\section{Integral morphisms and finite morphisms}
\label{section:II.6}

\subsection{Preschemes integral over another.}
\label{subsection:II.6.1}

\begin{definition}[6.1.1]
\label{II.6.1.1}
Let $X$ be an $S$-prescheme and $f:X\to S$ its structural morphism. We say
that $X$ is \emph{integral over $S$}, or that $f$ is an \emph{integral
morphism}, if there exists a covering $(S_\alpha)$ of $S$ by affine open
subsets such that, for every $\alpha$, the induced prescheme
$f^{-1}(S_\alpha)$ is an affine scheme whose ring $B_\alpha$ is an integral
algebra over the ring $A_\alpha$ of $S_\alpha$. We say that $X$ is
\emph{finite over $S$}, or that $f$ is a \emph{finite morphism}, if $X$ is
integral over $S$ and of finite type over $S$.
\end{definition}

When $S$ is affine with ring $A$, we will also say ``integral (resp. finite)
over $A$'' instead of ``integral (resp. finite) over $S$''.

\begin{env}[6.1.2]
\label{II.6.1.2}
It is clear that if $X$ is integral over $S$, it is affine over $S$. For a
prescheme $X$ affine over $S$ to be integral (resp. finite) over $S$, it is
necessary and sufficient that the associated quasi-coherent
$\mathcal{O}_S$-algebra $\mathcal{A}(X)$ be such that there exists a
covering $(S_\alpha)$ of $S$ by affine open subsets having the property that,
for every $\alpha$, $\Gamma(S_\alpha,\mathcal{A}(X))$ is an algebra integral
(resp. integral and of finite type) over
$\Gamma(S_\alpha,\mathcal{O}_S)$. A quasi-coherent $\mathcal{O}_S$-algebra
having this property is said to be \emph{integral} (resp. \emph{finite}) over
$\mathcal{O}_S$. To give a prescheme integral (resp. finite) over $S$
therefore amounts (\hyperref[II.1.3.1]{1.3.1}) to giving a
quasi-coherent $\mathcal{O}_S$-algebra integral (resp. finite) over
$\mathcal{O}_S$. Note that a quasi-coherent $\mathcal{O}_S$-algebra
$\mathcal{B}$ is finite if and only if it is an $\mathcal{O}_S$-module of
finite type (\hyperref[I.1.3.9]{I, 1.3.9}); it amounts to the same to say
that $\mathcal{B}$ is an $\mathcal{O}_S$-algebra that is integral and of
finite type, since an algebra that is integral and of finite type over a
ring $A$ is an $A$-module of finite type.
\end{env}

\begin{proposition}[6.1.3]
\label{II.6.1.3}
Let $S$ be a locally Noetherian prescheme. For a prescheme $X$ affine over
$S$ to be finite over $S$, it is necessary and sufficient that the
$\mathcal{O}_S$-algebra $\mathcal{A}(X)$ be coherent.
\end{proposition}

Taking into account the preceding remark, this amounts to observing that if
$S$ is locally Noetherian, the quasi-coherent $\mathcal{O}_S$-modules of
finite type are none other than the coherent $\mathcal{O}_S$-modules
(\hyperref[I.1.5.1]{I, 1.5.1}).

\begin{proposition}[6.1.4]
\label{II.6.1.4}
Let $X$ be a prescheme integral (resp. finite) over $S$ and $f:X\to S$ the
structural morphism. Then, for every affine open subset $U\subset S$, with
ring $A$, $f^{-1}(U)$ is an affine scheme whose ring $B$ is an integral
(resp. finite) algebra over $A$.
\end{proposition}

\oldpage[II]{111}

We will first prove the following lemma:

\begin{lemma}[6.1.4.1]
\label{II.6.1.4.1}
Let $A$ be a ring, $M$ an $A$-module, and $(g_i)_{1\leq i\leq m}$ a finite
system of elements of $A$ such that the $D(g_i)$ ($1\leq i\leq m$) cover
$\operatorname{Spec}(A)$. If, for every $i$, $M_{g_i}$ is an
$A_{g_i}$-module of finite type, then $M$ is an $A$-module of finite type.
\end{lemma}

Indeed, we may suppose that $M_{g_i}$ admits a finite system of generators
$(m_{ij}/g_i^n)$ with $m_{ij}\in M$, $n$ being the same for all the indices
$i$. Let us show that the $m_{ij}$ form a system of generators of $M$. Let
$M'$ be the sub-$A$-module of $M$ generated by these elements, and let $m$ be
an element of $M$. By hypothesis, for each $i$, there exist $a_{ij}\in A$ and
an integer $p$ (independent of $i$) such that, in $M_{g_i}$,
$m/1=(\sum_j a_{ij}m_{ij})/g_i^p$; this implies that there exists an integer
$r\geq p$ such that, for every $i$, we have $g_i^r m\in M'$. Now, since the
$D(g_i^r)=D(g_i)$ cover $\operatorname{Spec}(A)$, the ideal of $A$
generated by the $g_i^r$ is equal to $A$; in other words, there exist
elements $a_i\in A$ such that $\sum_i a_i g_i^r=1$; consequently
$m=(\sum_i a_i g_i^r)m\in M'$, whence the lemma.

That being so, we already know (\hyperref[II.1.3.2]{1.3.2}) that
$f^{-1}(U)$ is affine. If $\varphi$ is the homomorphism $A\to B$
corresponding to $f$, there exists a finite covering of $U$ by open subsets
$D(g_i)$ ($g_i\in A$) such that, if $h_i=\varphi(g_i)$, $B_{h_i}$ is an
integral (resp. integral and finite) algebra over $A_{g_i}$. Indeed, there
is a covering of $U$ by affine open subsets $V_\alpha\subset U$ such that, if
$A_\alpha=A(V_\alpha)$, $B_\alpha=A(f^{-1}(V_\alpha))$ is an integral (resp.
finite) algebra over $A_\alpha$. Every $x\in U$ belongs to some
$V_\alpha$, so there exists $g\in A$ such that $x\in D(g)\subset V_\alpha$;
if $g_\alpha$ is the image of $g$ in $A_\alpha$, we have
$A(D(g))=A_g=(A_\alpha)_{g_\alpha}$; let $h=\varphi(g)$, and let $h_\alpha$
be the image of $g_\alpha$ in $B_\alpha$; we have
\[
A(D(h))=B_h=(B_\alpha)_{h_\alpha}
\]
and since $B_\alpha$ is integral (resp. finite) over $A_\alpha$,
$(B_\alpha)_{h_\alpha}$ is integral (resp. finite) over
$(A_\alpha)_{g_\alpha}$. It now suffices to use the fact that $U$ is
quasi-compact to obtain the desired covering.

If we first suppose that the $B_{h_i}$ are integral and finite over the
$A_{g_i}$, then, since $B_{h_i}$, as an $A_{g_i}$-module, can also be
written $B_{g_i}$, Lemma (\hyperref[II.6.1.4.1]{6.1.4.1}) shows that in
this case $B$ is an $A$-module of finite type.

Now suppose only that each $B_{h_i}$ is integral over $A_{g_i}$; let
$b\in B$, and let $C$ be the sub-$A$-algebra of $B$ generated by $b$. For
every $i$, $C_{h_i}$ is the algebra over $A_{g_i}$ generated by $b/1$ in
$B_{h_i}$; it follows from the hypothesis that each $C_{h_i}$ is an
$A_{g_i}$-module of finite type, so
(\hyperref[II.6.1.4.1]{6.1.4.1}) $C$ is an $A$-module of finite type,
which proves that $B$ is integral over $A$.

\begin{proposition}[6.1.5]
\label{II.6.1.5}
\begin{enumerate}
\item[(i)] A closed immersion is finite (and \emph{a fortiori} integral).
\item[(ii)] The composite of two finite (resp. integral) morphisms is finite
(resp. integral).
\item[(iii)] If $f:X\to Y$ is a finite (resp. integral) $S$-morphism,
$f_{(S')}:X_{(S')}\to Y_{(S')}$ is finite (resp. integral) for every
extension $S'\to S$ of the base.
\item[(iv)] If $f:X\to Y$ and $g:X'\to Y'$ are two finite (resp. integral)
$S$-morphisms, then so is
$f\times_S g:X\times_S X'\to Y\times_S Y'$.
\item[(v)] If $f:X\to Y$ and $g:Y\to Z$ are two morphisms such that $g\circ
f$ is finite (resp. integral) and $g$ is separated, then $f$ is finite
(resp. integral).
\item[(vi)] If $f:X\to Y$ is a finite (resp. integral) morphism, then so is
$f_{\mathrm{red}}$.
\end{enumerate}
\end{proposition}

\oldpage[II]{112}

By virtue of (\hyperref[I.5.5.12]{I, 5.5.12}), it suffices to prove (i),
(ii) and (iii). To prove that a closed immersion $X\to S$ is finite, we may
restrict to the case where $S=\operatorname{Spec}(A)$, and everything then
comes down to observing that a quotient ring $A/\mathfrak{J}$ is a monogenic
$A$-module. To prove that the composite of two finite (resp. integral)
morphisms $X\to Y$, $Y\to Z$ is finite (resp. integral), we may again suppose
$Z$ (and consequently $X$, $Y$ (\hyperref[II.1.3.4]{1.3.4})) affine, and
then the assertion is equivalent to saying that if $B$ is a finite (resp.
integral) $A$-algebra and $C$ a finite (resp. integral) $B$-algebra, then
$C$ is a finite (resp. integral) $A$-algebra, which is immediate. Finally, to
prove (iii), we may restrict to the case where $S=Y$, since $X_{(S')}$ is
identified with $X\times_Y Y_{(S')}$ (\hyperref[I.3.3.11]{I, 3.3.11}); we
may moreover suppose that $S=\operatorname{Spec}(A)$,
$S'=\operatorname{Spec}(A')$; then $X$ is affine with ring $B$
(\hyperref[II.1.3.4]{1.3.4}), $X_{(S')}$ is affine with ring
$A'\otimes_A B$, and it suffices to observe that if $B$ is a finite (resp.
integral) $A$-algebra, $A'\otimes_A B$ is a finite (resp. integral)
$A'$-algebra.

Note also that if $X$ and $Y$ are two $S$-preschemes that are finite (resp.
integral) over $S$, their \emph{sum} $X\amalg Y$ is a prescheme finite
(resp. integral) over $S$, since this amounts to seeing that if $B$ and $C$
are two $A$-algebras finite (resp. integral) over $A$, then so is
$B\times C$.

\begin{corollary}[6.1.6]
\label{II.6.1.6}
If $X$ is a prescheme integral (resp. finite) over $S$, then, for every open
subset $U\subset S$, $f^{-1}(U)$ is integral (resp. finite) over $U$.
\end{corollary}

This is a particular case of (\hyperref[II.6.1.5]{6.1.5, (iii)}).

\begin{corollary}[6.1.7]
\label{II.6.1.7}
Let $f:X\to Y$ be a finite morphism. Then, for every $y\in Y$, the fibre
$f^{-1}(y)$ is an algebraic scheme finite over $k(y)$, and \emph{a fortiori}
its underlying space is discrete and finite.
\end{corollary}

Indeed, as a $k(y)$-prescheme, $f^{-1}(y)$ is identified with
$X\times_Y\operatorname{Spec}(k(y))$ (\hyperref[I.3.6.1]{I, 3.6.1}), which
is finite over $\operatorname{Spec}(k(y))$
(\hyperref[II.6.1.5]{6.1.5, (iii)}); it is therefore an affine scheme
whose ring is an algebra of finite rank over $k(y)$
(\hyperref[II.6.1.4]{6.1.4}). The proposition then follows from
(\hyperref[I.6.4.4]{I, 6.4.4}).

\begin{corollary}[6.1.8]
\label{II.6.1.8}
Let $X$, $S$ be two integral preschemes and $f:X\to S$ a dominant morphism.
If $f$ is integral (resp. finite), the field $R(X)$ of rational functions on
$X$ is algebraic (resp. algebraic of finite degree) over the field $R(S)$ of
rational functions on $S$.
\end{corollary}

Indeed, let $s$ be the generic point of $S$; the $k(s)$-prescheme
$f^{-1}(s)$ is integral (resp. finite) over
$\operatorname{Spec}(k(s))$ (\hyperref[II.6.1.5]{6.1.5, (iii)}), and by
hypothesis contains the generic point $x$ of $X$; the local ring of $x$ in
$f^{-1}(s)$, being equal to $k(x)$ (\hyperref[I.3.6.5]{I, 3.6.5}), is a
local ring of an algebra integral (resp. finite) over $k(s)$
(\hyperref[II.6.1.4]{6.1.4}), whence the corollary.

\begin{remark}[6.1.9]
\label{II.6.1.9}
The hypothesis that $g$ is \emph{separated} is essential for the validity of
(\hyperref[II.6.1.5]{6.1.5, (v)}): indeed, if $Y$ is not separated over
$Z$, the identity $1_Y$ is the composite morphism
$Y\xrightarrow{\Delta_Y}Y\times_ZY\xrightarrow{p_1}Y$, but $\Delta_Y$ is not
an integral morphism, as follows from
(\hyperref[II.6.1.10]{6.1.10}):
\end{remark}

\begin{proposition}[6.1.10]
\label{II.6.1.10}
Every integral morphism is universally closed.
\end{proposition}

Let $f:X\to Y$ be an integral morphism; by virtue of
(\hyperref[II.6.1.5]{6.1.5, (iii)}), it suffices to show that $f$ is
closed. Let $Z$ be a closed subset of $X$; there exists a sub-prescheme of
$X$ having

\oldpage[II]{113}

$Z$ as underlying space (\hyperref[I.5.2.1]{I, 5.2.1}), and it therefore
follows from (\hyperref[II.6.1.5]{6.1.5, (i) and (ii)}) that we may
restrict to proving that $f(X)$ is closed in $Y$. By virtue of
(\hyperref[II.6.1.5]{6.1.5, (vii)}), we may suppose $X$ and $Y$ reduced;
moreover, if $T$ is the reduced closed sub-prescheme of $Y$ having
$\overline{f(X)}$ as underlying space (\hyperref[I.5.2.1]{I, 5.2.1}), we
know that $f$ factors as $X\xrightarrow{g}T\xrightarrow{j}Y$, where $j$ is
the injection morphism (\hyperref[I.5.2.2]{I, 5.2.2}), and since $j$ is
separated (\hyperref[I.5.5.1]{I, 5.5.1, (i)}), it follows from
(\hyperref[II.6.1.5]{6.1.5, (v)}) that $g$ is an integral morphism. We may
therefore suppose that $f(X)$ is dense in $Y$. Finally, the question being
local on $Y$, we may restrict to the case where $Y=\operatorname{Spec}(A)$.
Then $X=\operatorname{Spec}(B)$, where $B$ is an $A$-algebra integral over
$A$ (\hyperref[II.6.1.4]{6.1.4}); moreover $A$ is reduced
(\hyperref[I.5.1.4]{I, 5.1.4}), and the hypothesis that $f(X)$ is dense in
$Y$ implies that the homomorphism $\varphi:A\to B$ corresponding to $f$ is
injective (\hyperref[I.1.2.7]{I, 1.2.7}). To say that under these
conditions $f(X)=Y$ means that every prime ideal of $A$ is the trace on $A$
of a prime ideal of $B$, which is none other than the first
Cohen--Seidenberg theorem (\cite[vol.~I, p.~257, Theorem~3]{I-13}).

\begin{corollary}[6.1.11]
\label{II.6.1.11}
Every finite morphism $f:X\to Y$ is projective.
\end{corollary}

Indeed, since $f$ is affine, $\mathcal{O}_X$ is an $\mathcal{O}_X$-module
very ample relative to $f$ (\hyperref[II.5.1.2]{5.1.2}); moreover
$f_*(\mathcal{O}_X)$ is a quasi-coherent $\mathcal{O}_Y$-module of finite
type (\hyperref[II.6.1.2]{6.1.2}); finally $f$ is separated, of finite
type and universally closed (\hyperref[II.6.1.10]{6.1.10}), and we are
therefore under the conditions of application of the criterion
(\hyperref[II.5.5.4]{5.5.4, (ii)}).

\begin{proposition}[6.1.12]
\label{II.6.1.12}
Let $f:X'\to X$ be a finite morphism, and let
$\mathcal{B}=f_*(\mathcal{O}_{X'})$ (which is a quasi-coherent
$\mathcal{O}_X$-algebra and an $\mathcal{O}_X$-module of finite type). Let
$\mathcal{F}'$ be a quasi-coherent $\mathcal{O}_{X'}$-module; for
$\mathcal{F}'$ to be locally free, of rank $r$, it is necessary and
sufficient that $f_*(\mathcal{F}')$ be a locally free $\mathcal{B}$-module of
rank $r$.
\end{proposition}

It is clear that if $f_*(\mathcal{F}')|U$ is isomorphic to
$\mathcal{B}^r|U$ ($U$ an open subset of $X$), then $\mathcal{F}'|f^{-1}(U)$
is isomorphic to $\mathcal{O}_{X'}^r|f^{-1}(U)$
(\hyperref[II.1.4.2]{1.4.2}). Conversely, suppose that $\mathcal{F}'$ is
locally free of rank $r$, and let us show that $f_*(\mathcal{F}')$ is locally
isomorphic to $\mathcal{B}^r$ as a $\mathcal{B}$-module. Let $x$ be a point
of $X$; as $U$ runs through a fundamental system of affine neighbourhoods of
$x$, $f^{-1}(U)$ runs through a fundamental system of affine neighbourhoods
(\hyperref[II.1.2.5]{1.2.5}) of the finite set $f^{-1}(x)$, since $f$ is
closed (\hyperref[II.6.1.10]{6.1.10}). The proposition then follows from
the following lemma:

\begin{lemma}[6.1.12.1]
\label{II.6.1.12.1}
Let $Y$ be a prescheme, $\mathcal{E}$ a locally free $\mathcal{O}_Y$-module
of rank $r$, and $Z$ a finite subset of $Y$ contained in an affine open
subset $V$. Then there exists a neighbourhood $U\subset V$ of $Z$ such that
$\mathcal{E}|U$ is isomorphic to $\mathcal{O}_Y^r|U$.
\end{lemma}

We may evidently suppose $Y$ affine; for every $z_i\in Z$, there exists in
the closure $\overline{\{z_i\}}$ at least one closed point $z'_i$
(\hyperref[0.2.1.3]{0, 2.1.3}); if $Z'$ is the set of the $z'_i$, every
neighbourhood of $Z'$ is a neighbourhood of $Z$, and we may therefore suppose
$Z$ closed in $Y$ and discrete. Consider the reduced closed sub-prescheme of
$Y$ having $Z$ as underlying space (\hyperref[I.5.2.1]{I, 5.2.1}), and let
$j:Z\to Y$ be the canonical injection; $j^*(\mathcal{E})=\mathcal{E}\otimes_Y
\mathcal{O}_Z$ is locally free of rank $r$ on the discrete scheme $Z$, and is
consequently isomorphic to $\mathcal{O}_Z^r$; in other words, there exist $r$
sections $s_i$ ($1\leq i\leq r$) of $\mathcal{E}\otimes_Y\mathcal{O}_Z$
over $Z$ such that the homomorphism
$\mathcal{O}_Z^r\to\mathcal{E}\otimes_Y\mathcal{O}_Z$ defined by these
sections is bijective. But $Y=\operatorname{Spec}(A)$ is affine, $Z$ is
defined by an ideal $\mathfrak{J}$ of $A$, and we have
$\mathcal{E}=\widetilde{M}$, where $M$ is an $A$-module; the $s_i$ are
elements

\oldpage[II]{114}

of $M\otimes_A(A/\mathfrak{J})$ and are therefore images of $r$ elements
$t_i\in M=\Gamma(Y,\mathcal{E})$. For every $z_j\in Z$ there is then a
neighbourhood $V_j$ of $z_j$ such that the restrictions of the $t_i$ to $V_j$
define an isomorphism
$\mathcal{O}_Y^r|V_j\to\mathcal{E}|V_j$
(\hyperref[0.5.5.4]{0, 5.5.4}); the neighbourhood $U$, union of the
$V_j$, therefore answers the question.

\begin{proposition}[6.1.13]
\label{II.6.1.13}
Let $g:X'\to X$ be an integral morphism of preschemes, $Y$ a normal locally
integral prescheme, and $f$ a rational map from $Y$ to $X'$ such that
$g\circ f$ is an everywhere defined rational map
(\hyperref[I.7.2.1]{I, 7.2.1}); then $f$ is everywhere defined.
\end{proposition}

If $f_1$ and $f_2$ are two morphisms (from dense open subsets of $Y$ to
$X'$) of the class $f$, it is clear that $g\circ f_1$ and $g\circ f_2$ are
equivalent morphisms, which justifies the notation $g\circ f$ for their
equivalence class. Recall also that when we suppose moreover that $Y$ is
locally Noetherian, the hypothesis that $Y$ is normal already implies that
$Y$ is locally integral (\hyperref[I.6.1.13]{I, 6.1.13}).

To prove (\hyperref[II.6.1.13]{6.1.13}), note first that the question is
local on $Y$, and we may consequently suppose that there exists in the class
of $g\circ f$ a morphism $h:Y\to X$. Consider the inverse image
$Y'=X'_{(h)}=X'_{(Y)}$, and note that the morphism
$g'=g_{(Y)}:Y'\to Y$ is integral
(\hyperref[II.6.1.5]{6.1.5, (iii)}). In view of the correspondence between
rational maps from $Y$ to $X'$ and rational $Y$-sections of $Y'$
(\hyperref[I.7.1.2]{I, 7.1.2}), we see that we are reduced to proving the
particular case of (\hyperref[II.6.1.13]{6.1.13}) where $X=Y$, in other
words the

\begin{corollary}[6.1.14]
\label{II.6.1.14}
Let $X$ be a normal locally integral prescheme, $g:X'\to X$ an integral
morphism, and $f$ a rational $X$-section of $X'$. Then $f$ is everywhere
defined.
\end{corollary}

The question being local on $X$, we may suppose $X$ integral, and $f$ is then
identified with a morphism from an open subset $U$ of $X$ to $X'$
(\hyperref[I.7.2.2]{I, 7.2.2}) which is a $U$-section of $g^{-1}(U)$.
Since $g$ is separated, $f$ is a closed immersion of $U$ into $g^{-1}(U)$
(\hyperref[I.5.4.6]{I, 5.4.6}); let $Z$ be the closed sub-prescheme of
$g^{-1}(U)$ associated to $f$ (\hyperref[I.4.2.1]{I, 4.2.1}), which is
isomorphic to $U$, hence integral; let $X_1$ be the reduced sub-prescheme of
$X'$ having as underlying space the closure $\overline{Z}$ of $Z$ in $X'$
(\hyperref[I.5.2.1]{I, 5.2.1}); $Z$ is then a sub-prescheme induced on an
open subset of $X_1$ (\hyperref[I.5.2.3]{I, 5.2.3}), and since it is
irreducible, so is $X_1$, which is consequently integral. The morphism $f$
can then be considered as a rational $X$-section of $X_1$; since the
restriction of $g$ to $X_1$ is an integral morphism
(\hyperref[II.6.1.5]{6.1.5, (i) and (ii)}), we are finally reduced to
proving (\hyperref[II.6.1.14]{6.1.14}) in the particular case where
$X'=X_1$, in other words, the

\begin{corollary}[6.1.15]
\label{II.6.1.15}
Let $X$ be a normal integral prescheme, $X'$ an integral prescheme, and
$g:X'\to X$ an integral morphism. If there exists a rational $X$-section $f$
of $X'$, then $g$ is an isomorphism.
\end{corollary}

The question being local on $X$, we may suppose $X$ affine with integral
ring $A$, and then $X'$ is affine with ring $A'$ integral over $A$
(\hyperref[II.6.1.4]{6.1.4}) and integral; moreover, the argument of
(\hyperref[II.6.1.14]{6.1.14}) shows that there is a dense open subset of
$X$ isomorphic to a dense open subset of $X'$, so $A$ and $A'$ have the same
field of fractions. On the other hand, by virtue of
(\hyperref[I.8.2.1.1]{I, 8.2.1.1}) and of the hypothesis that the
$\mathcal{O}_x$ are integrally closed, the ring $A$ is integrally closed,
so $A'=A$, which completes the proof of
(\hyperref[II.6.1.13]{6.1.13}).

\subsection{Quasi-finite morphisms.}
\label{subsection:II.6.2}

\begin{proposition}[6.2.1]
\label{II.6.2.1}
Let $f:X\to Y$ be a morphism locally of finite type and $x$ a point of $X$.
The following conditions are equivalent:

\oldpage[II]{115}

\begin{enumerate}
\item[a)] The point $x$ is isolated in its fibre $f^{-1}(f(x))$.
\item[b)] The ring $\mathcal{O}_x$ is a quasi-finite
$\mathcal{O}_{f(x)}$-module (\hyperref[0.7.4.1]{0, 7.4.1}).
\end{enumerate}
\end{proposition}

The question being evidently local on $X$ and on $Y$, we may suppose
$X=\operatorname{Spec}(A)$ and $Y=\operatorname{Spec}(B)$ affine, $A$ being
a $B$-algebra of finite type (\hyperref[I.6.3.3]{I, 6.3.3}). Moreover, we
can replace $X$ by $X\times_Y\operatorname{Spec}(\mathcal{O}_{f(x)})$
without changing either the fibre $f^{-1}(f(x))$ or the local ring
$\mathcal{O}_x$ (\hyperref[I.3.6.5]{I, 3.6.5}); we may therefore suppose
that $B$ is a local ring (equal to $\mathcal{O}_{f(x)}$); if $\mathfrak{n}$
is the maximal ideal of $B$, $f^{-1}(f(x))$ is then an affine scheme with
ring $A/\mathfrak{n}A$, of finite type over $k(f(x))=B/\mathfrak{n}$
(\hyperref[I.6.4.11]{I, 6.4.11}). That being so, if a) holds, we may
moreover suppose that $f^{-1}(f(x))$ consists of the point $x$ alone; hence
$A/\mathfrak{n}A$ is of finite rank over $B/\mathfrak{n}$
(\hyperref[I.6.4.4]{I, 6.4.4}), in other words $A$ is a quasi-finite
$B$-module. Conversely, if this is so, $f^{-1}(f(x))$ is an Artinian affine
scheme, hence discrete (\hyperref[I.6.4.4]{I, 6.4.4}); consequently $x$
is isolated in its fibre, and this shows that b) implies a).

\begin{corollary}[6.2.2]
\label{II.6.2.2}
Let $f:X\to Y$ be a morphism of finite type. The following conditions are
equivalent:
\begin{enumerate}
\item[a)] Every point $x\in X$ is isolated in its fibre $f^{-1}(f(x))$
(in other words, the subspace $f^{-1}(f(x))$ is discrete).
\item[b)] For every $x\in X$, the prescheme $f^{-1}(f(x))$ is a finite
$k(f(x))$-prescheme.
\item[c)] For every $x\in X$, the ring $\mathcal{O}_x$ is a quasi-finite
$\mathcal{O}_{f(x)}$-module.
\end{enumerate}
\end{corollary}

The equivalence of a) and c) follows from
(\hyperref[II.6.2.1]{6.2.1}). On the other hand, since $f^{-1}(f(x))$ is
an algebraic $k(f(x))$-prescheme (\hyperref[I.6.4.11]{I, 6.4.11}), the
equivalence of a) and b) follows from (\hyperref[I.6.4.4]{I, 6.4.4}).

\begin{definition}[6.2.3]
\label{II.6.2.3}
When $f:X\to Y$ is a morphism of finite type satisfying the equivalent
conditions of (\hyperref[II.6.2.2]{6.2.2}), we say that $f$ is
\emph{quasi-finite}, or that $X$ is quasi-finite over $Y$.
\end{definition}

It is clear that every finite morphism is quasi-finite
(\hyperref[II.6.1.8]{6.1.8}).

\begin{proposition}[6.2.4]
\label{II.6.2.4}
\begin{enumerate}
\item[(i)] An immersion $X\to Y$ which is closed, or such that $X$ is
Noetherian, is a quasi-finite morphism.
\item[(ii)] If $f:X\to Y$ and $g:Y\to Z$ are quasi-finite morphisms,
$g\circ f$ is quasi-finite.
\item[(iii)] If $X$ and $Y$ are $S$-preschemes and $f:X\to Y$ is a
quasi-finite $S$-morphism, then $f_{(S')}:X_{(S')}\to Y_{(S')}$ is
quasi-finite for every extension $g:S'\to S$ of the base.
\item[(iv)] If $f:X\to Y$ and $g:X'\to Y'$ are two quasi-finite
$S$-morphisms,
\[
f\times_S g:X\times_S Y\to X'\times_S Y'
\]
is quasi-finite.
\item[(v)] Let $f:X\to Y$ and $g:Y\to Z$ be two morphisms such that $g\circ
f$ is quasi-finite; then, if moreover $g$ is separated, or $X$ Noetherian, or
$X\times_ZY$ locally Noetherian, $f$ is quasi-finite.
\item[(vi)] If $f$ is quasi-finite, then so is $f_{\mathrm{red}}$.
\end{enumerate}
\end{proposition}

If $f:X\to Y$ is an immersion, every fibre consists of a single point, and
assertion (i) therefore follows from (\hyperref[I.6.3.4]{I, 6.3.4 (i) and
6.3.5}). To prove (ii), we first observe that $h=g\circ f$ is of finite type
(\hyperref[I.6.3.4]{I, 6.3.4 (ii)}); moreover, if $z=h(x)$ and
$y=f(x)$, $y$ is isolated in $g^{-1}(z)$, so there exists an open
neighbourhood $V$ of $y$ in $Y$ meeting no point of $g^{-1}(z)$ other than
$y$; $f^{-1}(V)$ is therefore an open neighbourhood of $x$ meeting none

\oldpage[II]{116}

of the $f^{-1}(y')$, where $y'\neq y$ is in $g^{-1}(z)$; since $x$ is
isolated in $f^{-1}(y)$, it is isolated in
$h^{-1}(z)=f^{-1}(g^{-1}(z))$. To prove (iii), we may restrict to the case
where $Y=S$ (\hyperref[I.3.3.11]{I, 3.3.11}); again we first observe that
$f'=f_{(S')}$ is of finite type
(\hyperref[I.6.3.4]{I, 6.3.4, (iii)}); on the other hand, if
$x'\in X'=X_{(S')}$ and if we set $y'=f'(x')$ and $y=g(y')$,
$f'^{-1}(y')$ is identified with
$f^{-1}(y)\otimes_{k(y)}k(y')$ (\hyperref[I.3.6.5]{I, 3.6.5}); since
$f^{-1}(y)$ is of finite rank over $k(y)$ by hypothesis, $f'^{-1}(y')$ is of
finite rank over $k(y')$, hence discrete. Assertions (iv), (v), (vi) are
deduced from (i), (ii), (iii) by the general method
(\hyperref[I.5.5.12]{I, 5.5.12}), except when the hypotheses in (v) are
other than the hypothesis ``$g$ separated''; to deal with this case, we
first observe that if $x$ is isolated in
$f^{-1}(g^{-1}(g(f(x))))$, it is \emph{a fortiori} isolated in
$f^{-1}(f(x))$; the fact that $f$ is of finite type then follows from
(\hyperref[I.6.3.6]{I, 6.3.6}).

\begin{proposition}[6.2.5]
\label{II.6.2.5}
Let $A$ be a complete Noetherian local ring, $X$ an $A$-scheme locally of
finite type, $x$ a point of $X$ above the closed point $y$ of
$Y=\operatorname{Spec}(A)$, and suppose $x$ isolated in its fibre
$f^{-1}(y)$ ($f$ being the structural morphism $X\to Y$). Then
$\mathcal{O}_x$ is an $A$-module of finite type and $X$ is $Y$-isomorphic to
the sum (I, 3.1) of
$X'=\operatorname{Spec}(\mathcal{O}_x)$ (which is a finite $Y$-scheme) and an
$A$-scheme $X''$.
\end{proposition}

It follows from (\hyperref[II.6.2.1]{6.2.1}) that $\mathcal{O}_x$ is a
quasi-finite $A$-module. Since $\mathcal{O}_x$ is Noetherian
(\hyperref[I.6.3.7]{I, 6.3.7}) and the homomorphism
$A\to\mathcal{O}_x$ is local, the hypothesis that $A$ is complete implies
that $\mathcal{O}_x$ is an $A$-module of finite type
(\hyperref[0.7.4.3]{0, 7.4.3}). Let
$X'=\operatorname{Spec}(\mathcal{O}_x)$ be the local scheme of $X$ at the
point $x$ (\hyperref[I.2.4.1]{I, 2.4.1}), and let $g:X'\to X$ be the
canonical morphism. Since the composite $X'\xrightarrow{g}X\xrightarrow{f}Y$
is finite (\hyperref[II.6.1.1]{6.1.1}) and $f$ is separated, $g$ is
finite (\hyperref[II.6.1.5]{6.1.5, (v)}), so $g(X')$ is closed in $X$
(\hyperref[II.6.1.10]{6.1.10}); on the other hand, since $g$ is of finite
type, $g$ is a local isomorphism at the closed point $x'$ of $X'$ by virtue
of the definition of $g$ and of (\hyperref[I.6.5.4]{I, 6.5.4}); but since
$X'$ is the only open neighbourhood of $x'$, this means that $g$ is an open
immersion, so $g(X')$ is also open in $X$, which completes the proof.

\begin{corollary}[6.2.6]
\label{II.6.2.6}
Let $A$ be a complete Noetherian local ring, $Y=\operatorname{Spec}(A)$, and
$f:X\to Y$ a separated and quasi-finite morphism. Then $X$ is $Y$-isomorphic
to a sum $X'\amalg X''$, where $X'$ is a finite $Y$-scheme and $X''$ a
quasi-finite $Y$-scheme such that, if $y$ is the closed point of $Y$,
$X''\cap f^{-1}(y)=\varnothing$.
\end{corollary}

Indeed, the fibre $f^{-1}(y)$ is finite and discrete by hypothesis, and the
corollary therefore follows, by induction on the number of points of this
fibre, from (\hyperref[II.6.2.5]{6.2.5}).

\begin{remark}[6.2.7]
\label{II.6.2.7}
In Chapter V, we will see that if $Y$ is locally Noetherian, a separated
quasi-finite morphism $X\to Y$ is necessarily quasi-affine.
\end{remark}

\subsection{Integral closure of a prescheme.}
\label{subsection:II.6.3}

\begin{proposition}[6.3.1]
\label{II.6.3.1}
Let $(X,\mathcal{A})$ be a ringed space, $\mathcal{B}$ a (commutative)
$\mathcal{A}$-algebra, and $f$ a section of $\mathcal{B}$ over $X$. The
following properties are equivalent:
\begin{enumerate}
\item[a)] The sub-$\mathcal{A}$-module of $\mathcal{B}$ generated by the
$f^n$ for $n\geq 0$ (\hyperref[0.5.1.1]{0, 5.1.1}) is of finite type.
\item[b)] There exists a sub-$\mathcal{A}$-algebra $\mathcal{C}$ of
$\mathcal{B}$ which is an $\mathcal{A}$-module of finite type and is such
that $f\in\Gamma(X,\mathcal{C})$.
\item[c)] For every $x\in X$, $f_x$ is integral over the stalk
$\mathcal{A}_x$.
\end{enumerate}
\end{proposition}

\oldpage[II]{117}

Since the sub-$\mathcal{A}$-module of $\mathcal{B}$ generated by the $f^n$ is
an $\mathcal{A}$-algebra, it is clear that a) implies b). On the other hand,
b) implies that for every $x\in X$, the $\mathcal{A}_x$-module
$\mathcal{C}_x$ is of finite type, which implies that every element of the
algebra $\mathcal{C}_x$, and in particular $f_x$, is integral over
$\mathcal{A}_x$. Finally, if for some $x\in X$ we have a relation of the form
\[
f_x^n+(a_1)_xf_x^{n-1}+\ldots+(a_n)_x=0
\]
the $a_i$ ($1\leq i\leq n$) being sections of $\mathcal{A}$ over an open
neighbourhood $U$ of $x$, the section
$f^n|U+a_1\cdot f^{n-1}|U+\ldots+a_n$ is zero over a neighbourhood
$V\subset U$ of $x$, from which it follows immediately that all the $f^k|V$
($k\geq0$) are linear combinations, with coefficients in
$\Gamma(V,\mathcal{A})$, of the $f^j|V$ for $0\leq j\leq n-1$; we conclude
that c) implies a).

When the equivalent conditions of (\hyperref[II.6.3.1]{6.3.1}) are
satisfied, we say that the section $f$ is \emph{integral over
$\mathcal{A}$}.

\begin{corollary}[6.3.2]
\label{II.6.3.2}
Under the hypotheses of (\hyperref[II.6.3.1]{6.3.1}), there exists a
(unique) sub-$\mathcal{A}$-module $\mathcal{A}'$ of $\mathcal{B}$ such that,
for every $x\in X$, $\mathcal{A}'_x$ is the set of germs
$f_x\in\mathcal{B}_x$ that are integral over $\mathcal{A}_x$. For every open
subset $U\subset X$, the sections of $\mathcal{A}'$ over $U$ are the sections
$f\in\Gamma(U,\mathcal{B})$ that are integral over
$\mathcal{A}|U$.
\end{corollary}

The existence of $\mathcal{A}'$ is immediate, by taking for
$\Gamma(U,\mathcal{A}')$ the set of $f\in\Gamma(U,\mathcal{B})$ such that
$f_x$ is integral over $\mathcal{A}_x$ for every $x\in U$. The second
assertion follows immediately from (\hyperref[II.6.3.1]{6.3.1}).

It is clear that $\mathcal{A}'$ is a sub-$\mathcal{A}$-algebra of
$\mathcal{B}$; we say that it is the \emph{integral closure of
$\mathcal{A}$ in $\mathcal{B}$}.

\begin{env}[6.3.3]
\label{II.6.3.3}
Let $(X,\mathcal{A})$, $(Y,\mathcal{B})$ be two ringed spaces, and
\[
g=(\psi,\theta):X\to Y
\]
a morphism. Let $\mathcal{C}$ (resp. $\mathcal{D}$) be an
$\mathcal{A}$-algebra (resp. a $\mathcal{B}$-algebra) and let
\[
u:\mathcal{D}\to\mathcal{C}
\]
be a $g$-morphism (\hyperref[0.4.4.1]{0, 4.4.1}). Then, if
$\mathcal{A}'$ (resp. $\mathcal{B}'$) is the integral closure of
$\mathcal{A}$ (resp. $\mathcal{B}$) in $\mathcal{C}$ (resp.
$\mathcal{D}$), the \emph{restriction} of $u$ to $\mathcal{B}'$ is a
$g$-morphism
\[
u':\mathcal{B}'\to\mathcal{A}'.
\]

Indeed, if $j$ is the canonical injection $\mathcal{B}'\to\mathcal{D}$, it
suffices to show that
\[
v=u^\sharp\circ g^*(j):g^*(\mathcal{B}')\to\mathcal{C}
\]
maps $g^*(\mathcal{B}')$ into $\mathcal{A}'$. Now, an element of
$(g^*(\mathcal{B}'))_x=\mathcal{B}'_{\psi(x)}\otimes_{\mathcal{B}_{\psi(x)}}
\mathcal{A}_x$ is integral over $\mathcal{A}_x$ by definition of
$\mathcal{B}'$, so the same is true of its image under $v_x$, which
establishes our assertion.
\end{env}

\begin{proposition}[6.3.4]
\label{II.6.3.4}
Let $X$ be a prescheme and $\mathcal{A}$ a quasi-coherent
$\mathcal{O}_X$-algebra. The integral closure $\mathcal{O}'_X$ of
$\mathcal{O}_X$ in $\mathcal{A}$ is then a quasi-coherent
$\mathcal{O}_X$-algebra, and for every affine open subset $U$ of $X$,
$\Gamma(U,\mathcal{O}'_X)$ is the integral closure of
$\Gamma(U,\mathcal{O}_X)$ in $\Gamma(U,\mathcal{A})$.
\end{proposition}

We may restrict to the case where $X=\operatorname{Spec}(B)$ is affine and
$\mathcal{A}=\widetilde{A}$, $A$ being a

\oldpage[II]{118}

$B$-algebra; let $B'$ be the integral closure of $B$ in $A$. Everything comes
down to seeing that, for every $x\in X$, an element of $A_x$ integral over
$B_x$ necessarily belongs to $B'_x$, which follows from the fact that, for a
commutative ring $C$, the operations of integral closure in a $C$-algebra and
of passage to a ring of fractions (for a multiplicative subset of $C$)
commute (\cite[vol.~I, pp.~261 and 257]{I-13}).

The $X$-scheme $X'=\operatorname{Spec}(\mathcal{O}'_X)$ is then called the
\emph{integral closure of $X$ relative to $\mathcal{A}$} (or \emph{relative
to $\operatorname{Spec}(\mathcal{A})$}); it is clear that $X'$ is
\emph{integral} over $X$ (\hyperref[II.6.1.2]{6.1.2}).

It follows immediately from (\hyperref[II.6.3.4]{6.3.4}) that if
$f:X'\to X$ is the structural morphism, then, for every open subset $U$ of
$X$, $f^{-1}(U)$ is the \emph{integral closure of the prescheme induced by
$X$ on $U$, relative to $\mathcal{A}|U$}.

\begin{env}[6.3.5]
\label{II.6.3.5}
Let $X$, $Y$ be two preschemes, $f:X\to Y$ a morphism, $\mathcal{A}$ (resp.
$\mathcal{B}$) a quasi-coherent $\mathcal{O}_X$-algebra (resp.
$\mathcal{O}_Y$-algebra), and let $u:\mathcal{B}\to\mathcal{A}$ be an
$f$-morphism. We have seen (\hyperref[II.6.3.3]{6.3.3}) that we deduce
from it an $f$-morphism $u':\mathcal{O}'_Y\to\mathcal{O}'_X$, where
$\mathcal{O}'_X$ (resp. $\mathcal{O}'_Y$) is the integral closure of
$\mathcal{O}_X$ (resp. $\mathcal{O}_Y$) in $\mathcal{A}$ (resp.
$\mathcal{B}$). Consequently, if $X'$ (resp. $Y'$) is the integral closure of
$X$ (resp. $Y$) relative to $\mathcal{A}$ (resp. $\mathcal{B}$), we
canonically deduce from $u$ a morphism $f'=\operatorname{Spec}(u'):X'\to Y'$
(\hyperref[II.1.5.6]{1.5.6}) making the diagram
\[
\label{II.6.3.5.1}
\xymatrix{
X'\ar[r]^{f'}\ar[d]&Y'\ar[d]\\
X\ar[r]_{f}&Y
}
\tag{6.3.5.1}
\]
commutative.
\end{env}

\begin{env}[6.3.6]
\label{II.6.3.6}
Suppose that $X$ has only a \emph{finite} number of irreducible components
$X_i$ ($1\leq i\leq r$), with generic points $\xi_i$, and consider in
particular the integral closure of $X$ relative to a quasi-coherent
$\mathcal{R}(X)$-algebra $\mathcal{A}$ (as an $\mathcal{O}_X$-algebra or as
an $\mathcal{R}(X)$-algebra, which amounts to the same). We know
(\hyperref[I.7.3.5]{I, 7.3.5}) that $\mathcal{A}$ is the direct product
of $r$ quasi-coherent $\mathcal{O}_X$-algebras $\mathcal{A}_i$, the support
of $\mathcal{A}_i$ being contained in $X_i$, and the sheaf induced by
$\mathcal{A}_i$ on $X_i$ being a simple sheaf whose stalk $A_i$ is an
algebra over $\mathcal{O}_{\xi_i}$. It is then clear
(\hyperref[II.6.3.4]{6.3.4}) that the integral closure $\mathcal{O}'_X$ of
$\mathcal{O}_X$ in $\mathcal{A}$ is the \emph{direct product} of the integral
closures $\mathcal{O}^{(i)}_X$ of $\mathcal{O}_X$ in each of the
$\mathcal{A}_i$, and consequently the integral closure
$X'=\operatorname{Spec}(\mathcal{O}'_X)$ of $X$ relative to $\mathcal{A}$ is
an $X$-scheme that is the \emph{sum} of the
$\operatorname{Spec}(\mathcal{O}^{(i)}_X)=X'_i$ ($1\leq i\leq r$).

Suppose moreover that the $\mathcal{O}_X$-algebra $\mathcal{A}$ is
\emph{reduced}, or, what amounts to the same, that each of the algebras $A_i$
is reduced, and can consequently be considered as an algebra over the
\emph{field} $k(\xi_i)$ (equal to the field of rational functions of the
reduced sub-prescheme of $X$ having $X_i$ as underlying space); then
(\hyperref[II.1.3.8]{1.3.8}) each of the $X'_i$ is a \emph{reduced}
$X$-scheme and $X'$ is also the integral closure of $X_{\mathrm{red}}$.
Suppose further that each of the algebras $A_i$ is the \emph{direct product
of finitely many fields} $K_{ij}$ ($1\leq j\leq s_i$); if
$\mathcal{K}_{ij}$ is the subalgebra of $\mathcal{A}_i$ corresponding to
$K_{ij}$, it is clear that $\mathcal{O}^{(i)}_X$ is the \emph{direct product}
of the integral closures $\mathcal{O}^{(ij)}_X$ of $\mathcal{O}_X$ in each of
the $\mathcal{K}_{ij}$. Consequently, $X'_i$ is then the \emph{sum} of the
$X$-schemes $X'_{ij}=\operatorname{Spec}(\mathcal{O}^{(ij)}_X)$
($1\leq j\leq s_i$). Moreover, under these hypotheses and with this
notation:
\end{env}

\oldpage[II]{119}

\begin{proposition}[6.3.7]
\label{II.6.3.7}
Each of the $X'_{ij}$ is an integral and normal $X$-prescheme, and its field
of rational functions $R(X'_{ij})$ is canonically identified with the
algebraic closure $K'_{ij}$ of $k(\xi_i)$ in $K_{ij}$.
\end{proposition}

By virtue of the foregoing, we may suppose $X$ integral, so $r=1$, and
$s_1=1$, so that the unique algebra $A_1$ is a field $K$; let $\xi$ be the
generic point of $X$, and let $f:X'\to X$ be the structural morphism. For
every non-empty affine open subset $U$ of $X$, $f^{-1}(U)$ is identified with
the spectrum of the integral closure $B'_U$, in the field $K$, of the
integral ring $B_U=\Gamma(U,\mathcal{O}_X)$ (\hyperref[II.6.3.4]{6.3.4});
since the ring $B'_U$ is integral and integrally closed, so are the local
rings of the points of its spectrum, so $f^{-1}(U)$ is by definition an
\emph{integral and normal} scheme ((\hyperref[0.4.1.4]{0, 4.1.4}) and
(\hyperref[I.5.1.4]{I, 5.1.4})). Moreover, since $(0)$ is the only prime
ideal of $B'_U$ above the prime ideal $(0)$ of $B_U$
(\cite[vol.~I, p.~259]{I-13}), $f^{-1}(\xi)$ consists of a single point
$\xi'$, and $k(\xi')$ is the field of fractions $K'$ of $B'_U$, which is none
other than the algebraic closure of $k(\xi)$ in $K$. Finally, $X'$ is
irreducible, since, as $U$ runs through the set of non-empty affine open
subsets of $X$, the $f^{-1}(U)$ form an open covering of $X'$ consisting of
irreducible open subsets; moreover, the intersection $f^{-1}(U\cap V)$ of
two of these open subsets contains $\xi'$, hence is not empty, and we
conclude by (\hyperref[0.2.1.4]{0, 2.1.4}).

\begin{corollary}[6.3.8]
\label{II.6.3.8}
Let $X$ be a reduced prescheme having only finitely many irreducible
components $X_i$ ($1\leq i\leq r$), and let $\xi_i$ be the generic point of
$X_i$. The integral closure $X'$ of $X$ relative to $\mathcal{R}(X)$ is the
sum of $r$ $X$-schemes $X'_i$ which are integral and normal. If $f:X'\to X$
is the structural morphism, $f^{-1}(\xi_i)$ consists only of the generic
point $\xi'_i$ of $X'_i$, and we have $k(\xi'_i)=k(\xi_i)$; in other words,
$f$ is birational.
\end{corollary}

We say in this case that $X'$ is the \emph{normalization} of the reduced
prescheme $X$; note that $f$, being birational and integral, is
\emph{surjective} (\hyperref[II.6.1.10]{6.1.10}). Then, for $X'=X$ to
hold, it is necessary and sufficient that $X$ be normal. When $X$ is an
integral prescheme, it follows from (\hyperref[II.6.3.8]{6.3.8}) that its
normalization $X'$ is integral.

\begin{env}[6.3.9]
\label{II.6.3.9}
Let $X$, $Y$ be two integral preschemes, $f:X\to Y$ a dominant morphism, and
$L=R(X)$, $K=R(Y)$ the fields of rational functions of $X$ and $Y$; to $f$
there canonically corresponds an injection $K\to L$, and when we identify
$K$ (resp. $L$) with the simple sheaf $\mathcal{R}(Y)$ (resp.
$\mathcal{R}(X)$), this injection is an $f$-morphism. Let $K_1$ (resp.
$L_1$) be an extension of $K$ (resp. $L$), and suppose given a monomorphism
$K_1\to L_1$ such that the diagram
\[
\xymatrix{
K_1\ar[r]&L_1\\
K\ar[u]\ar[r]&L\ar[u]
}
\]
is commutative; if $K_1$ (resp. $L_1$) is considered as a simple sheaf on $Y$
(resp. $X$), hence as an $\mathcal{R}(Y)$-algebra (resp. an
$\mathcal{R}(X)$-algebra), this means that $K_1\to L_1$ is an
$f$-morphism. That being so, if $X'$ (resp. $Y'$) is the integral closure of
$X$ (resp. $Y$) relative to $L_1$ (resp. $K_1$), $X'$ (resp. $Y'$) is an
integral and normal prescheme (\hyperref[II.6.3.6]{6.3.6}) whose field of
rational functions is canonically identified with the algebraic closure

\oldpage[II]{120}

$L'$ (resp. $K'$) of $L$ (resp. $K$) in $L_1$ (resp. $K_1$), and there exists
a canonical (necessarily dominant) morphism $f':X'\to Y'$ making the diagram
(\hyperref[II.6.3.5.1]{6.3.5.1}) commutative.

The most important case is the one where we take $L_1=L$, $K_1$ thus being
an extension of $K$ contained in $L$, and where we suppose $X$ integral and
\emph{normal}, so that $X'=X$. The foregoing then shows that, when $X$ is
normal and $Y'$ is the integral closure of $Y$ relative to a field
$K_1\subset L=R(X)$, every dominant morphism $f:X\to Y$ factors as
\[
f:X\xrightarrow{f'}Y'\to Y
\]
where $f'$ is dominant; moreover, when the monomorphism $K_1\to L$ is given,
$f'$ is necessarily unique, as we see by reducing to the case where $X$ and
$Y$ are affine. We therefore see that, for given $Y$, $L$ and a given
$K$-monomorphism $K_1\to L$, the integral closure $Y'$ of $Y$ relative to
$K_1$ is the solution of a \emph{universal problem}.
\end{env}

\begin{remark}[6.3.10]
\label{II.6.3.10}
Let us return to the hypotheses of (\hyperref[II.6.3.6]{6.3.6}),
supposing moreover that each of the algebras $A_i$ is of \emph{finite rank}
over $k(\xi_i)$ (which implies that $A_i$ is the direct product of finitely
many fields); we can assert in certain cases that the structural morphism
$X'\to X$ is not only integral, but even \emph{finite}. Let us restrict to
the case where $X$ is \emph{reduced}; the question being local on $X$, we may
moreover suppose that $X$ is affine with ring $C$, and that $C$ has only
finitely many minimal ideals $\mathfrak{p}_i$ ($1\leq i\leq r$), such that
the $C_i=C/\mathfrak{p}_i$ are integral; $X'$ will then be finite over $X$ if
the integral closure of each $C_i$ in every extension of finite degree of its
field of fractions is a $C$-module of finite type
(\hyperref[II.6.3.4]{6.3.4}). We know that this condition is always
satisfied when $C$ is an \emph{algebra of finite type over a field}
(\cite[vol.~I, p.~267, Theorem~9]{I-13}), or \emph{over $\mathbf{Z}$}
(\cite[vol.~I, p.~93, Theorem~3]{I-9}), or \emph{over a complete Noetherian
local ring} ([25], p.~298). We conclude that $X'\to X$ will be a finite
morphism when $X$ is a scheme of \emph{finite type over a field}, or over
$\mathbf{Z}$, or over a complete Noetherian local ring.
\end{remark}

\subsection{Determinant of an endomorphism of an $\mathcal{O}_X$-module.}
\label{subsection:II.6.4}

\begin{env}[6.4.1]
\label{II.6.4.1}
Let $A$ be a (commutative) ring, $E$ a free $A$-module of rank $n$, and $u$
an endomorphism of $E$; recall that, to define the \emph{characteristic
polynomial} of $u$, we consider the endomorphism $u\otimes1$ of the free
$A[T]$-module of rank $n$, $E\otimes_AA[T]$ ($T$ an indeterminate), and we
set
\[
\label{II.6.4.1.1}
P(u,T)=\det(T\cdot I-(u\otimes1))
\tag{6.4.1.1}
\]
($I$ the identity automorphism of $E\otimes_AA[T]$). We have
\[
\label{II.6.4.1.2}
P(u,T)=T^n-\sigma_1(u)T^{n-1}+\ldots+(-1)^n\sigma_n(u)
\tag{6.4.1.2}
\]
where $\sigma_i(u)$ is an element of $A$, equal to a homogeneous polynomial
of degree $i$ (with integer coefficients) in the entries of the matrix of
$u$ relative to an arbitrary basis of $E$; we say that the $\sigma_i(u)$ are
the \emph{elementary symmetric functions of $u$}, and we have in particular
$\sigma_1(u)=\operatorname{Tr}u$ and $\sigma_n(u)=\det u$. Recall that, by
virtue of the Hamilton--Cayley theorem, we have
\[
\label{II.6.4.1.3}
P(u,u)=u^n-\sigma_1(u)u^{n-1}+\ldots+(-1)^n\sigma_n(u)=0
\tag{6.4.1.3}
\]

\oldpage[II]{121}

which can also be written
\[
\label{II.6.4.1.4}
(\det u)\cdot I_E=uQ(u)
\tag{6.4.1.4}
\]
($I_E$ the identity automorphism of $E$), with
\[
\label{II.6.4.1.5}
Q(u)=(-1)^{n+1}(u^{n-1}-\sigma_1(u)u^{n-2}+\ldots+(-1)^{n-1}\sigma_{n-1}(u)).
\tag{6.4.1.5}
\]

Let $\varphi:A\to B$ be a ring homomorphism, making $B$ an $A$-algebra;
consider the $B$-module $E_{(B)}=E\otimes_AB$, which is free of rank $n$,
and the extension $u\otimes1$ of $u$ to an endomorphism of $E_{(B)}$; it is
immediate that we have $\sigma_i(u\otimes1)=\varphi(\sigma_i(u))$ for every
index $i$.
\end{env}

\begin{env}[6.4.2]
\label{II.6.4.2}
Suppose now that $A$ is an \emph{integral} ring, with field of fractions
$K$, and $E$ an $A$-module of \emph{finite type} (but not necessarily
free). Let $n$ be the \emph{rank} of $E$, that is, the rank of the free
$K$-module $E\otimes_AK$; to every endomorphism $u$ of $E$ there canonically
corresponds the endomorphism $u\otimes1$ of $E\otimes_AK$; by abuse of
language, we again call the polynomial $P(u\otimes1,T)$ the
\emph{characteristic polynomial} of $u$ and denote it by $P(u,T)$; its
coefficients $\sigma_i(u\otimes1)$ (which belong to $K$) will again be
denoted by $\sigma_i(u)$ and called the \emph{elementary symmetric functions
of $u$}; in particular $\det u=\det(u\otimes1)$ by definition. With this
notation, the formulas (\hyperref[II.6.4.1.3]{6.4.1.3}) to
(\hyperref[II.6.4.1.5]{6.4.1.5}) make sense and are still valid when we
interpret $u^j$ as the homomorphism $E\to E\otimes_AK$ that is the composite
of the endomorphism $u^j\otimes1=(u\otimes1)^j$ of $E\otimes_AK$ and the
canonical homomorphism $x\mapsto x\otimes1$.

If $F$ is the torsion submodule of $E$ and $E_0=E/F$, we have
$u(F)\subset F$, so, by passage to the quotients, $u$ gives an endomorphism
$u_0$ of $E_0$; moreover $E\otimes_AK$ is identified with $E_0\otimes_AK$ and
$u\otimes1$ with $u_0\otimes1$, so $\sigma_i(u)=\sigma_i(u_0)$ for
$1\leq i\leq n$.

When $E$ is \emph{torsion-free}, $E$ is canonically identified with a
sub-$A$-module of $E\otimes_AK$, and the relation $u\otimes1=0$ is
equivalent to $u=0$. When $E$ is a free $A$-module, the two definitions of
the $\sigma_i(u)$ given in (\hyperref[II.6.4.1]{6.4.1}) and
(\hyperref[II.6.4.2]{6.4.2}) coincide by the foregoing, which justifies
the notation adopted.

Note also that when $E$ is a \emph{torsion module}, in other words
$E_0=\{0\}$, the exterior algebra of $E_0$ reduces to $K$ and the
determinant of the unique endomorphism $u_0$ of $E_0$ is \emph{equal to
$1$}.
\end{env}

\begin{proposition}[6.4.3]
\label{II.6.4.3}
Let $A$ be an integral ring, $E$ an $A$-module of finite type, and $u$ an
endomorphism of $E$; then the elementary symmetric functions $\sigma_i(u)$
of $u$ (and in particular $\det u$) are elements of $K$ integral over $A$.
\end{proposition}

Let $A'$ be the integral closure of $A$; since $A'[T]$ is an integrally
closed ring ([24], p.~99), it is the integral closure of $A[T]$ in its field
of fractions $K(T)$. Replacing $u$ by $T\cdot I-u\otimes1$ and $A$ by $A[T]$,
we see that we are reduced to proving that $\det u$ is integral over $A$. If
$n$ is the rank of $E$, we have $\det u=\det(\bigwedge^n u)$ and
$(\bigwedge^n u)\otimes1=\bigwedge^n(u\otimes1)$, so we may suppose
$n=1$. But then the map $u\mapsto\det u$ is a homomorphism from the
$A$-module $\operatorname{Hom}_A(E,E)$ to $K$; since $E$ is of finite type,
$\operatorname{Hom}_A(E,E)$ is isomorphic to a submodule of the $A$-module of
finite type $E^n$ (if $E$ admits a system of $n$ generators); hence the
elements $\det u$ belong to a sub-$A$-module of $K$ of finite type, and are
consequently integral over $A$.

\oldpage[II]{122}

\begin{corollary}[6.4.4]
\label{II.6.4.4}
Under the hypotheses of (\hyperref[II.6.4.3]{6.4.3}), if we suppose
moreover that $A$ is normal, the $\sigma_i(u)$ (in particular
$\operatorname{Tr}u$ and $\det u$) belong to $A$.
\end{corollary}

\begin{proposition}[6.4.5]
\label{II.6.4.5}
Let $A$ be an integral ring, $E$ an $A$-module of finite type, of rank $n$,
and $u$ an endomorphism of $E$ such that the $\sigma_i(u)$ belong to $A$. For
$u$ to be an automorphism of $E$, it is necessary that $\det u$ be
invertible in $A$; this condition is sufficient when $E$ is torsion-free.
\end{proposition}

The condition is \emph{sufficient}, since the hypothesis and the formulas
(\hyperref[II.6.4.1.4]{6.4.1.4}) and
(\hyperref[II.6.4.1.5]{6.4.1.5}) (valid in $E$, and not only in
$E\otimes_AK$, since $E$ is torsion-free) prove that
$(\det u)^{-1}Q(u)$ is the inverse of $u$.

The condition is \emph{necessary}, since if $u$ is invertible, it follows
from (\hyperref[II.6.4.3]{6.4.3}) that $\det(u^{-1})$ belongs to the
integral closure $A'$ of $A$ in its field of fractions $K$ and is evidently
the inverse of $\det u$ in $A'$. Our assertion then follows from:

\begin{lemma}[6.4.5.1]
\label{II.6.4.5.1}
Let $A$ be a subring of a ring $A'$ such that $A'$ is integral over $A$. If
an element $x\in A$ is invertible in $A'$, it is also invertible in $A$.
\end{lemma}

Otherwise, $x$ would belong to a maximal ideal $\mathfrak{m}$ of $A$, and it
follows from the first Cohen--Seidenberg theorem
(\cite[vol.~I, p.~257, Theorem~3]{I-13}) that there would be a maximal ideal
$\mathfrak{m}'$ of $A'$ such that $\mathfrak{m}=\mathfrak{m}'\cap A$; we
would therefore have $x\in\mathfrak{m}'$, which is absurd.

\begin{corollary}[6.4.6]
\label{II.6.4.6}
Let $A$ be an integral and integrally closed ring, $E$ a torsion-free
$A$-module of finite type, and $u$ an endomorphism of $E$. For $u$ to be an
automorphism of $E$, it is necessary and sufficient that $\det u$ be
invertible in $A$.
\end{corollary}

This follows from (\hyperref[II.6.4.4]{6.4.4}) and
(\hyperref[II.6.4.5]{6.4.5}).

\begin{remark}[6.4.7]
\label{II.6.4.7}
Further on we will need a generalization of the preceding results. Consider
a \emph{reduced} Noetherian ring $A$; let $\mathfrak{p}_\alpha$
($1\leq\alpha\leq r$) be its minimal ideals, $K_\alpha$ the field of
fractions of the integral ring $A/\mathfrak{p}_\alpha$, and $K$ the total
ring of fractions of $A$, the \emph{direct product} of the fields
$K_\alpha$. Let $E$ be an $A$-module of finite type, and \emph{suppose that}
$E\otimes_AK$ is a free $K$-module of rank $n$ (which here is no longer a
consequence of the other hypotheses); it amounts to the same to say that all
the $K_\alpha$-vector spaces $E\otimes_AK_\alpha=E_\alpha$ have \emph{the
same dimension $n$}; if then $u$ is an endomorphism of $E$, we again set
$P(u,T)=P(u\otimes1,T)$ and $\sigma_j(u)=\sigma_j(u\otimes1)$, and in
particular $\det u=\det(u\otimes1)$; the $\sigma_j(u)$ are therefore
elements of $K$. It is immediate that $E\otimes_AK$ is the direct sum of the
$E_\alpha$ and that the latter are stable under $u\otimes1$; moreover, the
restriction of $u\otimes1$ to $E_\alpha$ is none other than the extension
$u_\alpha$ of $u$ to this $K_\alpha$-vector space; we conclude that
$\sigma_j(u)$ is the element of $K$ whose components in the $K_\alpha$ are
the $\sigma_j(u_\alpha)$. Since the integral closure of $A$ in $K$ is the
\emph{direct product} of the integral closures of $A$ in the $K_\alpha$, the
$\sigma_j(u)$ are \emph{integral} over $A$.

\begin{lemma}[6.4.7.1]
\label{II.6.4.7.1}
The sub-$A$-algebra of $K$ generated by all the elements $\sigma_j(u)$
($1\leq j\leq n$), as $u$ runs through $\operatorname{Hom}_A(E,E)$, is an
$A$-module of finite type.
\end{lemma}

It suffices to prove that the sub-$A[T]$-algebra of $K[T]$ generated by the
$P(u,T)$ is an $A[T]$-module of finite type, since, if the $F_i(T)$
($1\leq i\leq m$) form a system of generators of this $A[T]$-module, the
coefficients of the $F_i(T)$ are integral over $A$, and so generate an
$A$-algebra which is an $A$-module of finite type
(\cite[vol.~I, p.~255, Theorem~1]{I-13}). We can therefore

\oldpage[II]{123}

replace $A$ by $A[T]$ (which is Noetherian) and $E$ by
$E\otimes_AA[T]=E'$, which is such that
$E'\otimes_{A[T]}K[T]=E\otimes_AK[T]$ is a free $K[T]$-module of rank $n$.
Returning to the original notation, we therefore see that it suffices to
prove that the $A$-module generated by the elements $\det u$, where $u$ runs
through $\operatorname{Hom}_A(E,E)$, is of finite type; \emph{a fortiori}
(since every submodule of an $A$-module of finite type is of finite type) it
suffices to prove that, as $v$ runs through the set of endomorphisms of
$\bigwedge^nE$, the $A$-module generated by the $\det v$ is of finite type;
in other words, we are again reduced to the case where $n=1$. But then the
proposition follows from the fact that $\operatorname{Hom}_A(E,E)$ is an
$A$-module of finite type and that $v\mapsto\det v$ is a homomorphism from
this $A$-module to $K$.

Let $F$ be the kernel of the canonical homomorphism $E\to E\otimes_AK$ and
let $E_0=E/F$; we see as above that $E\otimes_AK$ is identified with
$E_0\otimes_AK$, that $u(F)\subset F$, and that, if $u_0$ is the
endomorphism of $E_0$ deduced from $u$ by passage to the quotients,
$u\otimes1$ is identified with $u_0\otimes1$, and
$\sigma_j(u)=\sigma_j(u_0)$ for every $j$. \emph{If $F=0$}, the formulas
(\hyperref[II.6.4.1.3]{6.4.1.3}) and
(\hyperref[II.6.4.1.5]{6.4.1.5}) make sense and are valid when $E$ is
identified with a submodule of $E\otimes_AK$ and the $u^j$ with
homomorphisms $E\to E\otimes_AK$; consequently, Proposition
(\hyperref[II.6.4.5]{6.4.5}) also extends to this case, together with its
proof.
\end{remark}

\begin{env}[6.4.8]
\label{II.6.4.8}
Let $(X,\mathcal{A})$ be a ringed space, $\mathcal{E}$ a \emph{locally free}
$\mathcal{A}$-module (of finite rank), and $u$ an endomorphism of
$\mathcal{E}$. By hypothesis there is a basis $\mathfrak{B}$ of the topology
of $X$ such that for every $V\in\mathfrak{B}$, $\mathcal{E}|V$ is isomorphic
to $\mathcal{A}^n|V$ (for an $n$ that may vary with $V$). Let $u$ be an
endomorphism of $\mathcal{E}$; for every $V\in\mathfrak{B}$, $u_V$ is
therefore an endomorphism of the $\Gamma(V,\mathcal{A})$-module
$\Gamma(V,\mathcal{E})$, which is free by hypothesis; the determinant
$\det u_V$ is therefore defined and belongs to $\Gamma(V,\mathcal{A})$.
Moreover, if $e_1,\ldots,e_n$ form a basis of $\Gamma(V,\mathcal{E})$, their
restrictions to every open subset $W\subset V$ form a basis of
$\Gamma(W,\mathcal{E})$ over $\Gamma(W,\mathcal{A})$, so $\det u_W$ is the
restriction of $\det u_V$ to $W$. There therefore exists one and only one
section of $\mathcal{A}$ over $X$, which we will denote by $\det u$ and call
the \emph{determinant of $u$}, such that the restriction of $\det u$ to every
$V\in\mathfrak{B}$ is $\det u_V$. It is clear that for every $x\in X$ we have
$(\det u)_x=\det u_x$; for two endomorphisms $u$, $v$ of $\mathcal{E}$, we
have
\[
\label{II.6.4.8.1}
\det(u\circ v)=(\det u)(\det v)
\tag{6.4.8.1}
\]
as well as
\[
\label{II.6.4.8.2}
\det(1_{\mathcal{E}})=1_{\mathcal{A}}
\tag{6.4.8.2}
\]
and, if the rank of $\mathcal{E}$ is constant (which will be the case
(\hyperref[0.5.4.1]{0, 5.4.1}) for $X$ connected) and equal to $n$,
\[
\label{II.6.4.8.3}
\det(s\cdot u)=s^n\det u
\tag{6.4.8.3}
\]
for every $s\in\Gamma(X,\mathcal{A})$ (note that $\det(0)=0_{\mathcal{A}}$
if $n\geq1$, but $\det(0)=1_{\mathcal{A}}$ for $n=0$). Moreover, for $u$ to
be an \emph{automorphism} of $\mathcal{E}$, it is necessary and sufficient
that $\det u$ be \emph{invertible} in $\Gamma(X,\mathcal{A})$.

If the rank of $\mathcal{E}$ is constant, we can define in the same way the
elementary symmetric functions $\sigma_i(u)$, which are elements of
$\Gamma(X,\mathcal{A})$; we again have the relations
(\hyperref[II.6.4.1.3]{6.4.1.3}) to
(\hyperref[II.6.4.1.5]{6.4.1.5}).

\oldpage[II]{124}

We have thus defined a homomorphism
$\det:\operatorname{Hom}_{\mathcal{A}}(\mathcal{E},\mathcal{E})
\to\Gamma(X,\mathcal{A})$ of multiplicative monoids; if we note that
$\operatorname{Hom}_{\mathcal{A}}(\mathcal{E},\mathcal{E})=
\Gamma(X,\mathcal{H}om_{\mathcal{A}}(\mathcal{E},\mathcal{E}))$ by
definition, we see that in this definition we can replace $X$ by an
arbitrary open subset $U$ of $X$, which immediately defines a homomorphism
$\det:\mathcal{H}om_{\mathcal{A}}(\mathcal{E},\mathcal{E})\to\mathcal{A}$ of
sheaves of multiplicative monoids. When $\mathcal{E}$ has constant rank, we
similarly define homomorphisms
$\sigma_i:\mathcal{H}om_{\mathcal{A}}(\mathcal{E},\mathcal{E})
\to\mathcal{A}$ of sheaves of sets; for $i=1$, the homomorphism
$\sigma_1=\operatorname{Tr}$ is a homomorphism of $\mathcal{A}$-modules.

Let $(Y,\mathcal{B})$ be a second ringed space, and let
$f:(X,\mathcal{A})\to(Y,\mathcal{B})$ be a morphism of ringed spaces; if
$\mathcal{F}$ is a locally free $\mathcal{B}$-module,
$f^*(\mathcal{F})$ is a locally free $\mathcal{A}$-module (which has the
same rank as $\mathcal{F}$ if the latter is of constant rank)
(\hyperref[0.5.4.5]{0, 5.4.5}). For every endomorphism $v$ of
$\mathcal{F}$, $f^*(v)$ is an endomorphism of $f^*(\mathcal{F})$, and it
follows immediately from the definitions that $\det f^*(v)$ is the section of
$\mathcal{A}=f^*(\mathcal{B})$ over $X$ which corresponds canonically to
$\det v\in\Gamma(Y,\mathcal{B})$. We can also say that the homomorphism
$f^*(\det):f^*(\mathcal{H}om_{\mathcal{B}}(\mathcal{F},\mathcal{F}))
\to f^*(\mathcal{B})=\mathcal{A}$ is the composite
\[
f^*(\mathcal{H}om_{\mathcal{B}}(\mathcal{F},\mathcal{F}))
\xrightarrow{\gamma^\sharp}
\mathcal{H}om_{\mathcal{A}}(f^*(\mathcal{F}),f^*(\mathcal{F}))
\xrightarrow{\det}\mathcal{A}
\]
(\hyperref[0.4.4.6]{0, 4.4.6}). We have analogous results for the
$\sigma_i$.
\end{env}

\begin{env}[6.4.9]
\label{II.6.4.9}
Suppose now that $X$ is a \emph{locally integral prescheme}, so that the
sheaf $\mathcal{R}(X)$ of rational functions on $X$ is a locally simple sheaf
of fields (\hyperref[I.7.3.4]{I, 7.3.4}), quasi-coherent as an
$\mathcal{O}_X$-module. If $\mathcal{E}$ is a quasi-coherent
$\mathcal{O}_X$-module \emph{of finite type},
$\mathcal{E}'=\mathcal{E}\otimes_{\mathcal{O}_X}\mathcal{R}(X)$ is then a
locally free $\mathcal{R}(X)$-module
(\hyperref[I.7.3.6]{I, 7.3.6}); for every endomorphism $u$ of
$\mathcal{E}$, $u\otimes1_{\mathcal{R}(X)}$ is then an endomorphism of
$\mathcal{E}'$, and $\det(u\otimes1)$ a section of $\mathcal{R}(X)$ over
$X$, which we again call the \emph{determinant of $u$} and again denote by
$\det u$. It follows from (\hyperref[II.6.4.3]{6.4.3}) that $\det u$ is
a section of the \emph{integral closure} of $\mathcal{O}_X$ in
$\mathcal{R}(X)$ (\hyperref[II.6.3.2]{6.3.2}); if moreover $X$ is
\emph{normal}, $\det u$ is therefore a \emph{section of $\mathcal{O}_X$}
over $X$, and if we suppose further that $\mathcal{E}$ is
\emph{torsion-free}, then for $u$ to be an automorphism of $\mathcal{E}$ it
is necessary and sufficient that $\det u$ be \emph{invertible}, by virtue of
(\hyperref[II.6.4.6]{6.4.6}). The formulas
(\hyperref[II.6.4.8.1]{6.4.8.1}) to
(\hyperref[II.6.4.8.3]{6.4.8.3}) are still valid; from the homomorphism
$u\mapsto\det u$, applied to the modules of sections of
$\mathcal{H}om_{\mathcal{O}_X}(\mathcal{E},\mathcal{E})$, we deduce a
homomorphism of sheaves
$\det:\mathcal{H}om_{\mathcal{O}_X}(\mathcal{E},\mathcal{E})
\to\mathcal{R}(X)$, which takes its values in $\mathcal{O}_X$ when $X$ is
normal. We have analogous definitions and results for the other elementary
symmetric functions $\sigma_j(u)$ when $\mathcal{E}'$ is of \emph{constant
rank}; if moreover $X$ is \emph{normal}, the $\sigma_j(u)$ are
\emph{sections of $\mathcal{O}_X$} over $X$.

Finally, let $X$, $Y$ be two integral preschemes, and let $f:X\to Y$ be a
\emph{dominant} morphism. We know that there then exists a canonical
homomorphism $f^*(\mathcal{R}(Y))\to\mathcal{R}(X)$
(\hyperref[I.7.3.8]{I, 7.3.8}\footnote{See \emph{Errata and Addenda},
list 1.}), whence, for every quasi-coherent $\mathcal{O}_Y$-module of
finite type $\mathcal{F}$, a canonical homomorphism
$\theta:f^*(\mathcal{F}\otimes_{\mathcal{O}_Y}\mathcal{R}(Y))
\to f^*(\mathcal{F})\otimes_{\mathcal{O}_X}\mathcal{R}(X)$. If $v$ is an
endo-

\oldpage[II]{125}

morphism of $\mathcal{F}$,
$f^*(v\otimes1_{\mathcal{R}(Y)})$ is an endomorphism of
$f^*(\mathcal{F}\otimes_{\mathcal{O}_Y}\mathcal{R}(Y))$, and we have a
commutative diagram
\[
\xymatrix{
f^*(\mathcal{F}\otimes_{\mathcal{O}_Y}\mathcal{R}(Y))
  \ar[r]^{f^*(v\otimes1)}\ar[d]_{\theta}
& f^*(\mathcal{F}\otimes_{\mathcal{O}_Y}\mathcal{R}(Y))\ar[d]^{\theta}\\
f^*(\mathcal{F})\otimes_{\mathcal{O}_X}\mathcal{R}(X)
  \ar[r]_{f^*(v)\otimes1}
& f^*(\mathcal{F})\otimes_{\mathcal{O}_X}\mathcal{R}(X).
}
\]
We easily conclude that $\det f^*(v)$ is the canonical image, under the
homomorphism $f^*(\mathcal{R}(Y))\to\mathcal{R}(X)$, of the section
$\det v$ of $\mathcal{R}(Y)$; indeed, we are immediately reduced to the case
where $X=\operatorname{Spec}(A)$ and $Y=\operatorname{Spec}(B)$ are affine,
$A$ and $B$ being integral rings with respective fields of fractions $K$ and
$L$, the homomorphism $B\to A$ being injective and therefore extending to a
monomorphism $L\to K$; if $\mathcal{F}=\widetilde{M}$, where $M$ is a
$B$-module of finite type, the rank over $L$ of $M\otimes_B L$ is
\emph{equal} to that of $(M\otimes_B A)\otimes_A K$ over $K$, and
$\det((u\otimes1)\otimes1)$ is the image in $K$ of $\det(u\otimes1)$ for
every endomorphism $u$ of $M$, whence the conclusion.
\end{env}

\begin{env}[6.4.10]
\label{II.6.4.10}
Suppose finally that $X$ is a \emph{locally Noetherian reduced prescheme},
so that the sheaf $\mathcal{R}(X)$ of rational functions on $X$ is again a
quasi-coherent $\mathcal{O}_X$-module
(\hyperref[I.7.3.4]{I, 7.3.4}); let, on the other hand, $\mathcal{E}$ be a
\emph{coherent} $\mathcal{O}_X$-module such that
$\mathcal{E}'=\mathcal{E}\otimes_{\mathcal{O}_X}\mathcal{R}(X)$ is
\emph{locally free} (of finite rank). By virtue of
(\hyperref[II.6.4.7]{6.4.7}), if $\mathcal{E}'$ is of constant rank, we
can, for every endomorphism $u$ of $\mathcal{E}$, define the $\sigma_j(u)$,
which are sections of $\mathcal{R}(X)$ over $X$. When we no longer suppose
$\mathcal{E}'$ of constant rank, we can still define the homomorphism
$\det:\mathcal{H}om_{\mathcal{O}_X}(\mathcal{E},\mathcal{E})
\to\mathcal{R}(X)$.
\end{env}

\subsection{Norm of an invertible sheaf.}
\label{subsection:II.6.5}

\begin{env}[6.5.1]
\label{II.6.5.1}
Let $(X,\mathcal{A})$ be a ringed space and let $\mathcal{B}$ be a
(commutative) $\mathcal{A}$-algebra. The $\mathcal{A}$-module $\mathcal{B}$
is canonically identified with a sub-$\mathcal{A}$-module of
$\mathcal{H}om_{\mathcal{A}}(\mathcal{B},\mathcal{B})$, a section $f$ of
$\mathcal{B}$ over an open subset $U$ of $X$ being identified with
multiplication by this section. If $(X,\mathcal{A})$ and $\mathcal{B}$
satisfy one of the conditions listed in
(\hyperref[II.6.4.8]{6.4.8}), (\hyperref[II.6.4.9]{6.4.9}) or
(\hyperref[II.6.4.10]{6.4.10}), we can therefore define $\det(f)$ (and in
certain cases the $\sigma_j(f)$), which are sections of $\mathcal{A}$ or of
$\mathcal{R}(X)$ over $U$, which we will call the \emph{norm of $f$} (resp.
the \emph{elementary symmetric functions of $f$}) and denote by
$N_{\mathcal{B}/\mathcal{A}}(f)$. We will suppose that \emph{one} of the
following conditions holds:
\begin{enumerate}
\item[(I)] \emph{$\mathcal{B}$ is a locally free $\mathcal{A}$-module
(of finite rank).}
\item[(II)] \emph{$(X,\mathcal{A})$ is a locally Noetherian reduced
prescheme, $\mathcal{B}$ is a coherent $\mathcal{A}$-module such that
$\mathcal{B}\otimes_{\mathcal{A}}\mathcal{R}(X)$ is a locally free
$\mathcal{R}(X)$-module, and for every section
$f\in\Gamma(U,\mathcal{B})$ over an open subset $U\subset X$,
$N_{\mathcal{B}/\mathcal{A}}(f)$ is a section of $\mathcal{A}$ over $U$.}
\end{enumerate}

\oldpage[II]{126}

Note that this last condition is automatically satisfied when the locally
Noetherian prescheme $X$ is \emph{normal}
(\hyperref[II.6.4.9]{6.4.9}).

On the other hand, the hypothesis that
$\mathcal{B}\otimes_{\mathcal{A}}\mathcal{R}(X)$ is locally free can also be
expressed as follows: denote by $X_\alpha$ the reduced closed sub-preschemes
of $X$ having as underlying spaces the irreducible components of $X$
(\hyperref[I.5.2.1]{I, 5.2.1}), which are therefore locally Noetherian
integral preschemes. Every $x\in X$ belongs to finitely many of the
subspaces $X_\alpha$; on the other hand,
$\mathcal{B}\otimes_{\mathcal{A}}\mathcal{R}(X_\alpha)$ is a locally free
$\mathcal{R}(X_\alpha)$-module of constant rank $k_\alpha$
(\hyperref[I.7.3.6]{I, 7.3.6}); to say that
$\mathcal{B}\otimes_{\mathcal{A}}\mathcal{R}(X)$ is a locally free
$\mathcal{R}(X)$-module means that, for every $x\in X$, \emph{the ranks
$k_\alpha$ corresponding to the indices such that $x\in X_\alpha$ are
equal}. Indeed, the question is local, and we are reduced to the case
$X=\operatorname{Spec}(C)$, where $C$ is a reduced Noetherian ring, and
$\mathcal{B}=\widetilde{D}$, where $D$ is a $C$-algebra which is a
$C$-module of finite type; if $\mathfrak{p}_i$ ($1\leq i\leq m$) are the
minimal prime ideals of $C$, the total ring of fractions $L$ of $C$ is the
direct product of the fields of fractions $K_i$ of the integral rings
$C/\mathfrak{p}_i$, and $D\otimes_C L$ is the direct sum of the
$D\otimes_C K_i$, whence the conclusion.

It is clear that under the hypotheses (I) or (II) we have thus defined a
\emph{homomorphism of sheaves of multiplicative monoids}
$N_{\mathcal{B}/\mathcal{A}}:\mathcal{B}\to\mathcal{A}$, also denoted by $N$
if no confusion results, and called the \emph{norm} homomorphism. For two
sections $f$, $g$ of $\mathcal{B}$ over the same open subset $U$, we
therefore have
\[
\label{II.6.5.1.1}
N_{\mathcal{B}/\mathcal{A}}(fg)
=N_{\mathcal{B}/\mathcal{A}}(f)N_{\mathcal{B}/\mathcal{A}}(g)
\tag{6.5.1.1}
\]
for the corresponding sections of $\mathcal{A}$ over $U$;
\[
\label{II.6.5.1.2}
N_{\mathcal{B}/\mathcal{A}}(1_{\mathcal{B}})=1_{\mathcal{A}}~;
\tag{6.5.1.2}
\]
finally, for every section $s$ of $\mathcal{A}$ over $U$,
\[
\label{II.6.5.1.3}
N_{\mathcal{B}/\mathcal{A}}(s\cdot1_{\mathcal{B}})=s^n
\tag{6.5.1.3}
\]
if the rank of $\mathcal{B}$ is constant and equal to $n$ (for
$s=0_{\mathcal{A}}$, this formula gives
$N(0_{\mathcal{B}})=0_{\mathcal{A}}$ if $n\geq1$,
$N(0_{\mathcal{B}})=N(1_{\mathcal{B}})=1_{\mathcal{A}}$ if $n=0$).

Under hypothesis (I), for $f\in\Gamma(U,\mathcal{B})$ to be invertible, it
is necessary and sufficient that $N(f)\in\Gamma(U,\mathcal{A})$ be so. Under
hypothesis (II), this condition is necessary; it is sufficient (by
(\hyperref[II.6.4.7]{6.4.7})) when we suppose that
$\mathcal{B}\to\mathcal{B}\otimes_{\mathcal{A}}\mathcal{R}(X)$ is
\emph{injective} and that the following more restrictive hypothesis holds:
\begin{enumerate}
\item[(II \emph{bis})] \emph{$(X,\mathcal{A})$ is a locally Noetherian
reduced prescheme, $\mathcal{B}$ is a coherent $\mathcal{A}$-module such that
$\mathcal{B}\otimes_{\mathcal{A}}\mathcal{R}(X)$ is a locally free
$\mathcal{R}(X)$-module, and for every section
$f\in\Gamma(U,\mathcal{B})$ over an open subset such that
$\mathcal{B}\otimes_{\mathcal{A}}\mathcal{R}(X)|U$ is of constant rank $n$
over $\mathcal{R}(X)|U$, the $\sigma_j(f)$ ($1\leq j\leq n$) are sections of
$\mathcal{A}$ over $U$.}
\end{enumerate}
(Note again that this condition is satisfied when $X$ is
\emph{normal}.)
\end{env}

\begin{env}[6.5.2]
\label{II.6.5.2}
Suppose that one of the hypotheses (I), (II) of
(\hyperref[II.6.5.1]{6.5.1}) holds, and let $\mathcal{L}'$ be an
\emph{invertible} $\mathcal{B}$-module. We will associate to it canonically
(up to unique isomorphism) an \emph{invertible} $\mathcal{A}$-module
$\mathcal{L}$ in the following way. Denote by $\mathcal{A}^*$ (resp.
$\mathcal{B}^*$) the sub-sheaf of $\mathcal{A}$ (resp. $\mathcal{B}$) such
that $\Gamma(U,\mathcal{A}^*)$ (resp. $\Gamma(U,\mathcal{B}^*)$) is the set
of

\oldpage[II]{127}

invertible elements of $\Gamma(U,\mathcal{A})$ (resp.
$\Gamma(U,\mathcal{B})$) for every open subset $U\subset X$; these are
sheaves of multiplicative groups, and $N_{\mathcal{B}/\mathcal{A}}$,
restricted to $\mathcal{B}^*$, is a \emph{homomorphism}
$\mathcal{B}^*\to\mathcal{A}^*$ of sheaves of groups
(\hyperref[II.6.5.1]{6.5.1}). Let $\mathfrak{L}$ be the set of pairs
$(U_\lambda,\eta_\lambda)$ having the following property: $U_\lambda$ is an
open subset of $X$, and $\eta_\lambda$ is an isomorphism
$\mathcal{L}'|U_\lambda\simeq\mathcal{B}|U_\lambda$ of
$(\mathcal{B}|U_\lambda)$-modules. By hypothesis, the $U_\lambda$ form a
covering of $X$; for any two indices $\lambda$, $\mu$, we set
$\omega_{\lambda\mu}=(\eta_\lambda|U_\lambda\cap U_\mu)
\circ(\eta_\mu|U_\lambda\cap U_\mu)^{-1}$, an automorphism of
$\mathcal{B}|U_\lambda\cap U_\mu$, which is canonically identified with a
section of $\mathcal{B}^*$ over $U_\lambda\cap U_\mu$, and
$(\omega_{\lambda\mu})$ is a $1$-cocycle of the covering
$\mathfrak{U}=(U_\lambda)$ with values in $\mathcal{B}^*$
(\hyperref[0.5.4.7]{0, 5.4.7}). The fact that
$N_{\mathcal{B}/\mathcal{A}}:\mathcal{B}^*\to\mathcal{A}^*$ is a
homomorphism then implies that
$(N_{\mathcal{B}/\mathcal{A}}\omega_{\lambda\mu})$ is a $1$-cocycle of
$\mathfrak{U}$ with values in $\mathcal{A}^*$, and there therefore
corresponds to it (up to unique isomorphism) an invertible
$\mathcal{A}$-module $\mathcal{L}$ (\hyperref[0.5.4.7]{0, 5.4.7}), which
we will denote by $N_{\mathcal{B}/\mathcal{A}}(\mathcal{L}')$ and call
\emph{the norm of the invertible $\mathcal{B}$-module $\mathcal{L}'$}.

Let $\mathfrak{M}$ be a subset of $\mathfrak{L}$ such that the corresponding
$U_\lambda$ still form a covering of $X$, and let $\mathfrak{B}$ be this
covering; the restriction of the cocycle $(\omega_{\lambda\mu})$ to
$\mathfrak{B}$ defines a $1$-cocycle
$(N_{\mathcal{B}/\mathcal{A}}\omega_{\lambda\mu})$ of $\mathfrak{B}$ with
values in $\mathcal{A}^*$, the restriction of the $1$-cocycle
$(N_{\mathcal{B}/\mathcal{A}}\omega_{\lambda\mu})$ of $\mathfrak{U}$; it is
clear that there is a canonical isomorphism of the invertible
$\mathcal{A}$-module defined by this $1$-cocycle of $\mathfrak{B}$ onto
$N_{\mathcal{B}/\mathcal{A}}(\mathcal{L}')$, which makes it possible to
define this invertible $\mathcal{A}$-module by an arbitrary subcovering of
$\mathfrak{U}$. This possibility immediately shows that, if
$\mathcal{L}'_1$, $\mathcal{L}'_2$ are two invertible $\mathcal{B}$-modules,
we have, by virtue of (\hyperref[II.6.5.1.1]{6.5.1.1}) and
(\hyperref[II.6.5.1.2]{6.5.1.2}),
\[
\label{II.6.5.2.1}
N(\mathcal{L}'_1\otimes_{\mathcal{B}}\mathcal{L}'_2)
=N(\mathcal{L}'_1)\otimes_{\mathcal{A}}N(\mathcal{L}'_2)
\tag{6.5.2.1}
\]
and
\[
\label{II.6.5.2.2}
N_{\mathcal{B}/\mathcal{A}}(\mathcal{B})=\mathcal{A}
\tag{6.5.2.2}
\]
as well as
\[
\label{II.6.5.2.3}
N((\mathcal{L}')^{-1})=(N(\mathcal{L}'))^{-1}
\tag{6.5.2.3}
\]
up to canonical isomorphisms. Moreover, it follows from
(\hyperref[II.6.5.1.3]{6.5.1.3}) that if $\mathcal{L}$ is an invertible
$\mathcal{A}$-module and $\mathcal{B}$ is of constant rank $n$ over
$\mathcal{A}$ in case (I) (resp.
$\mathcal{B}\otimes_{\mathcal{A}}\mathcal{R}(X)$ is of constant rank $n$
over $\mathcal{R}(X)$ in case (II)), we have, up to canonical isomorphism,
\[
\label{II.6.5.2.4}
N_{\mathcal{B}/\mathcal{A}}(\mathcal{L}\otimes_{\mathcal{A}}\mathcal{B})
=\mathcal{L}^{\otimes n}.
\tag{6.5.2.4}
\]
\end{env}

\begin{env}[6.5.3]
\label{II.6.5.3}
Let us now show that $N_{\mathcal{B}/\mathcal{A}}(\mathcal{L}')$ is a
\emph{covariant functor} on the category of invertible
$\mathcal{B}$-modules. Indeed, let $h':\mathcal{L}'_1\to\mathcal{L}'_2$ be a
homomorphism of invertible $\mathcal{B}$-modules, and let
$\mathfrak{B}=(U_\lambda)$ be an open covering of $X$ such that for every
$\lambda$ we have $(\mathcal{B}|U_\lambda)$-isomorphisms
$\eta^{(1)}_\lambda:\mathcal{L}'_1|U_\lambda\simeq\mathcal{B}|U_\lambda$ and
$\eta^{(2)}_\lambda:\mathcal{L}'_2|U_\lambda\simeq\mathcal{B}|U_\lambda$;
there is therefore for each $\lambda$ an endomorphism $h'_\lambda$ of
$\mathcal{B}|U_\lambda$ such that
$h'_\lambda\circ\eta^{(1)}_\lambda
=\eta^{(2)}_\lambda\circ(h'|U_\lambda)$, and we can evidently identify
$h'_\lambda$ with a section of $\mathcal{B}$ over $U_\lambda$
(\hyperref[0.5.1.1]{0, 5.1.1}). Hence, for every pair $(\lambda,\mu)$ of
indices, the restrictions to $U_\lambda\cap U_\mu$ of
$(\eta^{(2)}_\lambda)^{-1}\circ h'_\lambda\circ\eta^{(1)}_\lambda$ and of
$(\eta^{(2)}_\mu)^{-1}\circ h'_\mu\circ\eta^{(1)}_\mu$ coincide. We deduce
from this, for the $1$-cocycles $(\omega^{(1)}_{\lambda\mu})$ and
$(\omega^{(2)}_{\lambda\mu})$ with values in $\mathcal{B}^*$ corresponding to
$\mathcal{L}'_1$ and $\mathcal{L}'_2$, the relation
\[
\omega^{(2)}_{\lambda\mu}h'_\mu
=h'_\lambda\omega^{(1)}_{\lambda\mu}.
\]

\oldpage[II]{128}

If we set $h_\lambda=N(h'_\lambda)$, we will therefore have the analogous
relations
\[
N(\omega^{(2)}_{\lambda\mu})h_\mu
=h_\lambda N(\omega^{(1)}_{\lambda\mu})
\]
and consequently the $h_\lambda$ define a homomorphism
$N(\mathcal{L}'_1)\to N(\mathcal{L}'_2)$, which we will denote by
$N_{\mathcal{B}/\mathcal{A}}(h')$ or $N(h')$. Under hypothesis (I), for
$h'$ to be an \emph{isomorphism}, it is necessary and sufficient (since this
is a local question) that $N_{\mathcal{B}/\mathcal{A}}(h')$ be one. Under
hypothesis (II), this condition is still necessary; it is sufficient if
hypothesis (II \emph{bis}) holds and
$\mathcal{B}\to\mathcal{B}\otimes_{\mathcal{A}}\mathcal{R}(X)$ is injective.

Take in particular $\mathcal{L}'_1=\mathcal{B}$; the homomorphisms
$\mathcal{B}\to\mathcal{L}'$ are then identified
(\hyperref[0.5.1.1]{0, 5.1.1}) with the sections of $\mathcal{L}'$ over
$X$, whence a canonical map
\[
N_{\mathcal{B}/\mathcal{A}}:\Gamma(X,\mathcal{L}')
\to\Gamma(X,N_{\mathcal{B}/\mathcal{A}}(\mathcal{L}')).
\]
It again follows from (\hyperref[II.6.5.1.1]{6.5.1.1}) that if
$f'_1\in\Gamma(X,\mathcal{L}'_1)$, $f'_2\in\Gamma(X,\mathcal{L}'_2)$, we
have
\[
\label{II.6.5.3.1}
N(f'_1\otimes f'_2)=N(f'_1)\otimes N(f'_2).
\tag{6.5.3.1}
\]
For every invertible $\mathcal{A}$-module $\mathcal{L}$ and every section
$f\in\Gamma(X,\mathcal{L})$, we have, taking into account
(\hyperref[II.6.5.2.4]{6.5.2.4}),
\[
\label{II.6.5.3.2}
N_{\mathcal{B}/\mathcal{A}}(f\otimes1_{\mathcal{B}})=f^{\otimes n}
\tag{6.5.3.2}
\]
when $\mathcal{B}$ is of constant rank $n$ under hypothesis (I) (resp. when
$\mathcal{B}\otimes_{\mathcal{A}}\mathcal{R}(X)$ is of constant rank $n$
under hypothesis (II)). Finally, for the homomorphism
$\mathcal{B}\to\mathcal{L}'$ corresponding to a section $f'$ of
$\mathcal{L}'$ over $X$ to be an isomorphism, it is necessary and sufficient
that $f'_x$ be a basis of $\mathcal{L}'_x$ for every $x\in X$; under
hypothesis (I), this condition is therefore equivalent to saying that
$(N(f'))_x$ is a basis of $(N(\mathcal{L}'))_x$ for every $x$; under
hypothesis (II), this condition is still necessary, and it is sufficient
when $\mathcal{B}$ satisfies (II \emph{bis}) and
$\mathcal{B}\to\mathcal{B}\otimes_{\mathcal{A}}\mathcal{R}(X)$ is injective.
\end{env}

\begin{env}[6.5.4]
\label{II.6.5.4}
Let $(X,\mathcal{A})$, $(X',\mathcal{A}')$ be two ringed spaces,
$f:X'\to X$ a morphism, $\mathcal{B}$ an $\mathcal{A}$-algebra, and
$\mathcal{B}'=f^*(\mathcal{B})$ the inverse image $\mathcal{A}'$-algebra.
Suppose that one of the following hypotheses holds:
\begin{enumerate}
\item[$1^\circ$] $\mathcal{B}$ satisfies hypothesis (I) of
(\hyperref[II.6.5.1]{6.5.1}).
\item[$2^\circ$] $(X,\mathcal{A})$ and $\mathcal{B}$ satisfy hypothesis
(II) of (\hyperref[II.6.5.1]{6.5.1}), $(X',\mathcal{A}')$ is a locally
Noetherian reduced prescheme, and, if we denote by $X_\alpha$ and
$X'_\beta$ the reduced closed sub-preschemes of $X$ and $X'$ respectively
having as underlying spaces the irreducible components of these spaces, the
restriction of $f$ to each $X'_\beta$ is a \emph{dominant} morphism from
$X'_\beta$ to some $X_\alpha$.
\end{enumerate}
Under these conditions, $\mathcal{B}'$ satisfies hypothesis (I) or
hypothesis (II) of (\hyperref[II.6.5.1]{6.5.1}) respectively; the first
assertion is immediate. To establish the second, it suffices to see that for
every $x'\in X'$, the ranks of the
$\mathcal{B}'\otimes_{\mathcal{O}_{X'}}\mathcal{R}(X'_\beta)$ for the $\beta$
such that $x'\in X'_\beta$ are \emph{the same}. Now, if the restriction of
$f$ to $X'_\beta$ is a dominant morphism into $X_\alpha$, the rank of
$\mathcal{B}'\otimes_{\mathcal{O}_{X'}}\mathcal{R}(X'_\beta)$ is equal to that
of $\mathcal{B}\otimes_{\mathcal{O}_X}\mathcal{R}(X_\alpha)$ (as we see
immediately by reducing to the affine case as in
(\hyperref[II.6.4.9]{6.4.9})), whence our assertion by virtue of
hypothesis (II) and of (\hyperref[II.6.5.1]{6.5.1}).

\oldpage[II]{129}

That being so, it follows from (\hyperref[II.6.4.8]{6.4.8}),
(\hyperref[II.6.4.9]{6.4.9}) and (\hyperref[II.6.4.10]{6.4.10}) that if
$s$ is a section of $\mathcal{B}$ over an open subset $U\subset X$ and
$s'$ the corresponding section of $\mathcal{B}'$ over $f^{-1}(U)$, then
$N_{\mathcal{B}'/\mathcal{A}'}(s')$ is the section of $\mathcal{A}'$ over
$f^{-1}(U)$ which corresponds to the section
$N_{\mathcal{B}/\mathcal{A}}(s)$ of $\mathcal{A}$ over $U$.

If $\mathcal{M}$ is an invertible $\mathcal{B}$-module, we deduce from the
foregoing that, if $\mathcal{M}'=f^*(\mathcal{M})$ (which is an invertible
$\mathcal{B}'$-module), we have
$N_{\mathcal{B}'/\mathcal{A}'}(\mathcal{M}')
=f^*(N_{\mathcal{B}/\mathcal{A}}(\mathcal{M}))$ up to canonical
isomorphism.
\end{env}

\begin{env}[6.5.5]
\label{II.6.5.5}
Suppose from now on that $(X,\mathcal{A})$ is a \emph{prescheme}. Giving a
quasi-coherent $\mathcal{A}$-algebra $\mathcal{B}$ which is an
$\mathcal{A}$-module \emph{of finite type} is then equivalent, as we know,
to giving a \emph{finite} morphism $g:X'\to X$ such that
$g_*(\mathcal{O}_{X'})=\mathcal{B}$, defined up to $X$-isomorphism
(\hyperref[II.6.1.2]{6.1.2} and \hyperref[II.1.3.1]{1.3.1}). Moreover,
giving a quasi-coherent $\mathcal{O}_{X'}$-module $\mathcal{F}'$ is
equivalent to giving a quasi-coherent $\mathcal{B}$-module $\mathcal{F}$ such
that $g_*(\mathcal{F}')=\mathcal{F}$ (\hyperref[II.1.4.3]{1.4.3}), and for
$\mathcal{F}'$ to be invertible, it is necessary and sufficient that
$\mathcal{F}$ be so (\hyperref[II.6.1.12]{6.1.12}). To translate the
preceding results in terms of the finite morphism $g$, we will therefore have
to suppose either that $g_*(\mathcal{O}_{X'})$ is a \emph{locally free}
$\mathcal{O}_X$-module (of finite type), or that $(X,\mathcal{O}_X)$ and
$g_*(\mathcal{O}_{X'})$ satisfy hypothesis (II). For every invertible
$\mathcal{O}_{X'}$-module $\mathcal{L}'$, we will then set
\[
\label{II.6.5.5.1}
N_{X'/X}(\mathcal{L}')
=N_{g_*(\mathcal{O}_{X'})/\mathcal{O}_X}(g_*(\mathcal{L}'))
\tag{6.5.5.1}
\]
and we will say that it is the \emph{norm} (relative to $g$) of
$\mathcal{L}'$. Similarly, if $h':\mathcal{L}'_1\to\mathcal{L}'_2$ is a
homomorphism of invertible $\mathcal{O}_{X'}$-modules, we will set
\[
\label{II.6.5.5.2}
N_{X'/X}(h')
=N_{g_*(\mathcal{O}_{X'})/\mathcal{O}_X}(g_*(h')):
N_{X'/X}(\mathcal{L}'_1)\to N_{X'/X}(\mathcal{L}'_2).
\tag{6.5.5.2}
\]
In particular, for $\mathcal{L}'_1=\mathcal{O}_{X'}$, we thus obtain a
canonical map
\[
\label{II.6.5.5.3}
N_{X'/X}:\Gamma(X',\mathcal{L}')
\to\Gamma(X,N_{X'/X}(\mathcal{L}')).
\tag{6.5.5.3}
\]
We will leave most of these translations to the reader, and will restrict
ourselves to making the following ones explicit:
\end{env}

\begin{proposition}[6.5.6]
\label{II.6.5.6}
Let $g:X'\to X$ be a finite morphism, and suppose either that
$g_*(\mathcal{O}_{X'})$ is a locally free $\mathcal{O}_X$-module, or that
$(X,\mathcal{O}_X)$ and $g_*(\mathcal{O}_{X'})$ satisfy (II \emph{bis})
(which will be the case in particular when $X$ is locally Noetherian and
normal). For a homomorphism $h':\mathcal{L}'_1\to\mathcal{L}'_2$ of
invertible $\mathcal{O}_{X'}$-modules to be an isomorphism, it is necessary
and sufficient, under the first hypothesis, that $N_{X'/X}(h')$ be an
isomorphism; under the second hypothesis, this condition is still necessary,
and it is sufficient when the homomorphism
$g_*(\mathcal{O}_{X'})\to
g_*(\mathcal{O}_{X'})\otimes_{\mathcal{O}_X}\mathcal{R}(X)$ is injective.
\end{proposition}

Note that we use here the fact that, for $g_*(h')$ to be an isomorphism, it
is necessary and sufficient that $h'$ be one
(\hyperref[II.1.4.2]{1.4.2}).

\begin{corollary}[6.5.7]
\label{II.6.5.7}
Let $g:X'\to X$ be a finite morphism, and suppose either that
$g_*(\mathcal{O}_{X'})$ is a locally free $\mathcal{O}_X$-module, or that
$(X,\mathcal{O}_X)$ and $g_*(\mathcal{O}_{X'})$ satisfy (II \emph{bis})
and that
$g_*(\mathcal{O}_{X'})\to
g_*(\mathcal{O}_{X'})\otimes_{\mathcal{O}_X}\mathcal{R}(X)$ is injective.
Let $\mathcal{L}'$ be an invertible $\mathcal{O}_{X'}$-module, $f'$ a
section of $\mathcal{L}'$ over $X'$, and $f=N_{X'/X}(f')$ the section of
$\mathcal{L}=N_{X'/X}(\mathcal{L}')$ over $X$ which corresponds to it
(\hyperref[II.6.5.5.3]{6.5.5.3}).

\oldpage[II]{130}

Then we have $g(X'-X'_{f'})=X-X_f$, and $X_f$ is the largest open subset
$U$ of $X$ such that $g^{-1}(U)\subset X'_{f'}$.
\end{corollary}

Indeed, $g(X'-X'_{f'})$ is closed in $X$
(\hyperref[II.6.1.10]{6.1.10}); it therefore suffices to prove the last
assertion. Now the relation $U\subset X_f$ is equivalent to the fact that the
homomorphism $\mathcal{O}_X|U\to\mathcal{L}|U$ defined by $f|U$ is an
isomorphism. By virtue of (\hyperref[II.6.5.6]{6.5.6}), this is
equivalent to saying that the homomorphism
$\mathcal{O}_{X'}|g^{-1}(U)\to\mathcal{L}'|g^{-1}(U)$ defined by
$f'|g^{-1}(U)$ is an isomorphism, that is, to the relation
$g^{-1}(U)\subset X'_{f'}$.

\begin{proposition}[6.5.8]
\label{II.6.5.8}
Let $g:X'\to X$ be a finite morphism and $f:Y\to X$ a morphism; let
$Y'=X'_{(Y)}$, $g'=g_{(Y)}$, $f'=f_{(X')}$, so that we have the commutative
diagram
\[
\xymatrix{
X'\ar[d]_g & Y'\ar[l]_{f'}\ar[d]^{g'}\\
X & Y\ar[l]^f
}
\]
Suppose either that $g_*(\mathcal{O}_{X'})$ is locally free, or that
$(X,\mathcal{O}_X)$ and $g_*(\mathcal{O}_{X'})$ satisfy (II), that $Y$ is a
locally Noetherian reduced prescheme, and that the restriction of $f$ to
every irreducible component of $Y$ is a dominant morphism into an irreducible
component of $X$. Then, for every invertible $\mathcal{O}_{X'}$-module
$\mathcal{L}'$, we have
\[
N_{Y'/Y}({f'}^*(\mathcal{L}'))=f^*(N_{X'/X}(\mathcal{L}'))
\]
up to canonical isomorphism.
\end{proposition}

Note that we have $f^*(g_*(\mathcal{L}'))=g'_*({f'}^*(\mathcal{L}'))$ by
virtue of (\hyperref[II.1.5.2]{1.5.2}), and in particular
$g'_*(\mathcal{O}_{Y'})=f^*(g_*(\mathcal{O}_{X'}))$; if
$g_*(\mathcal{O}_{X'})$ is locally free, then so is
$g'_*(\mathcal{O}_{Y'})$. The conclusion then follows from the definitions
and from (\hyperref[II.6.5.4]{6.5.4}).

\begin{remark}[6.5.9]
\label{II.6.5.9}
We will later generalize the notion of norm developed above, and relate it
to the notion of the direct image of a divisor.
\end{remark}

\subsection{Application: ampleness criteria.}
\label{subsection:II.6.6}

\begin{proposition}[6.6.1]
\label{II.6.6.1}
Let $Y$ be a prescheme, $f:X\to Y$ a quasi-compact morphism, and
$g:X'\to X$ a finite and surjective morphism. Suppose either that
$g_*(\mathcal{O}_{X'})$ is a locally free $\mathcal{O}_X$-module, or
that $(X,\mathcal{O}_X)$ and $g_*(\mathcal{O}_{X'})$ satisfy (II
\emph{bis}). Then, for every invertible $\mathcal{O}_{X'}$-module
$\mathcal{L}'$ ample for $f\circ g$, $N_{X'/X}(\mathcal{L}')=\mathcal{L}$
is ample for $f$.
\end{proposition}

We may suppose $Y$ affine (\hyperref[II.4.6.4]{4.6.4}), and then, by
virtue of (\hyperref[II.4.6.6]{4.6.6}), the statement is equivalent to

\begin{corollary}[6.6.2]
\label{II.6.6.2}
Let $X$ be a quasi-compact prescheme and $g:X'\to X$ a finite surjective
morphism such that $g_*(\mathcal{O}_{X'})$ is a locally free
$\mathcal{O}_X$-module, or such that $(X,\mathcal{O}_X)$ and
$g_*(\mathcal{O}_{X'})$ satisfy (II \emph{bis}). Then, for every ample
$\mathcal{O}_{X'}$-module $\mathcal{L}'$,
$\mathcal{L}=N_{X'/X}(\mathcal{L}')$ is ample.
\end{corollary}

Under the second hypothesis, we may moreover suppose that the canonical
homomorphism
$g_*(\mathcal{O}_{X'})\to
g_*(\mathcal{O}_{X'})\otimes_{\mathcal{O}_X}\mathcal{R}(X)$ is
\emph{injective}. Indeed, otherwise, let $\mathcal{T}$

\oldpage[II]{131}

be the kernel of this homomorphism, which is a coherent ideal of
$g_*(\mathcal{O}_{X'})=\mathcal{B}$
(\hyperref[I.6.1.1]{I, 6.1.1}), and set
$X''=\operatorname{Spec}(\mathcal{B}/\mathcal{T})$; we thus have a
commutative diagram
\[
\xymatrix{
X''\ar[rr]^h\ar[dr]_{g'} && X'\ar[dl]^g\\
& X
}
\]
where $h$ is a closed immersion (\hyperref[II.1.4.10]{1.4.10}). Moreover,
we know that the support of $\mathcal{T}$ is a closed set
(\hyperref[0.5.2.2]{0, 5.2.2}) which is rare in $X$
(\hyperref[I.7.4.6]{I, 7.4.6}), from which we conclude that for the generic
point $x$ of an irreducible component of $X$, there is an affine open
neighbourhood $U$ of $x$ such that
$\mathcal{B}|U=(\mathcal{B}/\mathcal{T})|U$. Since $g$ is surjective by
hypothesis, we conclude that $x\in g'(X'')$; $g'$ is therefore dominant, and,
being a finite morphism, it is \emph{surjective}
(\hyperref[II.6.1.10]{6.1.10}); by definition we have
\[
g'_*(\mathcal{O}_{X''})\otimes_{\mathcal{O}_X}\mathcal{R}(X)
=(\mathcal{B}/\mathcal{T})\otimes_{\mathcal{O}_X}\mathcal{R}(X)
=g_*(\mathcal{O}_{X'})\otimes_{\mathcal{O}_X}\mathcal{R}(X),
\]
hence $(X,\mathcal{O}_X)$ and $g'_*(\mathcal{O}_{X''})$ satisfy
(II \emph{bis}), and moreover
\[
g'_*(\mathcal{O}_{X''})\longrightarrow
g'_*(\mathcal{O}_{X''})\otimes_{\mathcal{O}_X}\mathcal{R}(X)
\]
is injective. Finally, $h^*(\mathcal{L}')=\mathcal{L}''$ is an ample
$\mathcal{O}_{X''}$-module
(\hyperref[II.4.6.13]{4.6.13, (i \emph{bis})}), and we have
$N_{X''/X}(\mathcal{L}'')=N_{X'/X}(\mathcal{L}')$. Indeed, to define
these two invertible $\mathcal{O}_X$-modules we can use the same
affine open covering $(U_\lambda)$ of $X$ such that the restrictions of
$g_*(\mathcal{L}')$ and $g'_*(\mathcal{L}'')$ to $U_\lambda$ are
isomorphic to $\mathcal{B}|U_\lambda$ and
$(\mathcal{B}/\mathcal{T})|U_\lambda$ respectively. We see immediately that to
every isomorphism
\[
\eta_\lambda:g_*(\mathcal{L}')|U_\lambda\longrightarrow
\mathcal{B}|U_\lambda
\]
there corresponds canonically an isomorphism
\[
\eta'_\lambda:g'_*(\mathcal{L}'')|U_\lambda\longrightarrow
(\mathcal{B}/\mathcal{T})|U_\lambda,
\]
in such a way that, if $(\omega_{\lambda\mu})$ and
$(\omega'_{\lambda\mu})$ are the $1$-cocycles corresponding to the systems
of isomorphisms $(\eta_\lambda)$ and $(\eta'_\lambda)$
(\hyperref[II.6.5.2]{6.5.2}), $\omega'_{\lambda\mu}$ is the canonical
image in
$\Gamma(U_\lambda\cap U_\mu,\mathcal{B}/\mathcal{T})$ of
$\omega_{\lambda\mu}\in\Gamma(U_\lambda\cap U_\mu,\mathcal{B})$. By virtue
of the definition of $\mathcal{T}$, we conclude that
\[
N_{\mathcal{B}/\mathcal{A}}\omega_{\lambda\mu}
=N_{(\mathcal{B}/\mathcal{T})/\mathcal{A}}\omega'_{\lambda\mu}
\]
(with $\mathcal{A}=\mathcal{O}_X$), whence the stated equality.

Suppose therefore that the homomorphism
$g_*(\mathcal{O}_{X'})\to
g_*(\mathcal{O}_{X'})\otimes_{\mathcal{O}_X}\mathcal{R}(X)$ is injective
when we are under hypothesis (II \emph{bis}). It suffices to prove that,
as $f$ runs through the sections of $\mathcal{L}^{\otimes n}$ ($n>0$)
over $X$, the $X_f$ form a basis of the topology of $X$
(\hyperref[II.4.5.2]{4.5.2}). Now let $x\in X$, and let $U$ be an arbitrary
neighbourhood of $x$; since $g^{-1}(x)$ is finite
(\hyperref[II.6.1.7]{6.1.7}) and $\mathcal{L}'$ is ample, there exist an
integer $n>0$ and a section $f'$ of ${\mathcal{L}'}^{\otimes n}$ over
$X'$ such that $X'_{f'}$ is a neighbourhood of $g^{-1}(x)$ contained in
$g^{-1}(U)$ (\hyperref[II.4.5.4]{4.5.4}). Since
\[
\mathcal{L}^{\otimes n}=N_{X'/X}({\mathcal{L}'}^{\otimes n}),
\]
it suffices to take $f=N_{X'/X}(f')$; indeed, we then have
$X-X_f=g(X'-X'_{f'})$ (\hyperref[II.6.5.7]{6.5.7}), hence
$x\in X_f\subset U$.

\begin{corollary}[6.6.3]
\label{II.6.6.3}
Under the hypotheses of (\hyperref[II.6.6.1]{6.6.1}), for an invertible
$\mathcal{O}_X$-module $\mathcal{L}$ to be ample for $f$, it is necessary
and sufficient that $\mathcal{L}'=g^*(\mathcal{L})$ be ample for $f\circ g$.
\end{corollary}

The condition is necessary, since $g$ is affine
(\hyperref[II.5.1.12]{5.1.12}). To prove that the condition is
sufficient, we may suppose $Y$ affine (\hyperref[II.4.6.4]{4.6.4}), hence
$X$ and $X'$ quasi-compact and $\mathcal{L}'$ ample
(\hyperref[II.4.6.6]{4.6.6}), and we must show that $\mathcal{L}$ is
ample. Now the set of points

\oldpage[II]{132}

$x\in X$ in a neighbourhood of which $g_*(\mathcal{O}_{X'})$ (resp.
$g_*(\mathcal{O}_{X'})\otimes_{\mathcal{O}_X}\mathcal{R}(X)$) has a given
rank $n$ under the first (resp. the second) hypothesis is both open
and closed in $X$; hence $X$ is the sum prescheme of finitely many of these
open subsets, and we may consequently suppose that it is equal to one of them
(\hyperref[II.4.6.17]{4.6.17}). But then we have
$N_{X'/X}(\mathcal{L}')=\mathcal{L}^{\otimes n}$, hence
$\mathcal{L}^{\otimes n}$ is ample by virtue of
(\hyperref[II.6.6.2]{6.6.2}), and therefore so is $\mathcal{L}$
(\hyperref[II.4.5.6]{4.5.6}).

\begin{corollary}[6.6.4]
\label{II.6.6.4}
Suppose that the hypotheses of
(\hyperref[II.6.6.1]{6.6.1}) are satisfied, and suppose moreover that
$f:X\to Y$ is of finite type. Then, for $f$ to be quasi-projective, it is
necessary and sufficient that $f\circ g$ be so. If we suppose in addition
that $Y$ is a quasi-compact scheme or a prescheme whose underlying space is
Noetherian, then, for $f$ to be projective, it is necessary and sufficient
that $f\circ g$ be so.
\end{corollary}

The hypothesis implies that $f\circ g$ is of finite type. Taking into account
the definition of quasi-projective morphisms
(\hyperref[II.5.3.1]{5.3.1}), the first assertion follows from
(\hyperref[II.6.6.1]{6.6.1}) and
(\hyperref[II.6.6.3]{6.6.3}). Taking into account this result and
(\hyperref[II.5.5.3]{5.5.3, (ii)}), it remains to prove that, when $f$ is
supposed quasi-projective, for $f$ to be proper it is necessary and
sufficient that $f\circ g$ be so. Now $f$ is then separated
(\hyperref[II.5.3.1]{5.3.1}) and of finite type; since $g$ is surjective,
our assertion follows from (\hyperref[II.5.4.2]{5.4.2, (ii)}) and
(\hyperref[II.5.4.3]{5.4.3, (ii)}).

In particular:

\begin{corollary}[6.6.5]
\label{II.6.6.5}
Let $X$ be a prescheme of finite type over a field $K$ and $K'$ an extension
of finite degree of $K$. For $X$ to be projective (resp. quasi-projective)
over $K$, it is necessary and sufficient that $X'=X\otimes_K K'$ be
projective (resp. quasi-projective) over $K'$.
\end{corollary}

The condition is indeed necessary
(\hyperref[II.5.3.4]{5.3.4, (iii)} and
\hyperref[II.5.5.5]{5.5.5, (iii)}). Conversely, suppose that it is
satisfied, and let $g:X'\to X$ be the canonical projection. It is clear that
$g$ is a finite morphism (\hyperref[II.6.1.5]{6.1.5, (iii)}) and that it is
surjective (\hyperref[I.3.5.2]{I, 3.5.2, (ii)}). Moreover,
$g_*(\mathcal{O}_{X'})$ is a locally free $\mathcal{O}_X$-module, being
isomorphic to $\mathcal{O}_X\otimes_K K'$
(\hyperref[II.1.5.2]{1.5.2}). It then follows from the hypothesis and from
(\hyperref[II.6.1.11]{6.1.11}) and
(\hyperref[II.5.5.5]{5.5.5, (ii)}) that $X'$ is projective (resp.
quasi-projective) \emph{over $K$}; we then deduce from
(\hyperref[II.6.6.4]{6.6.4}) that $X$ is projective (resp.
quasi-projective) over $K$.

In Chapter V, we will show that the statement
(\hyperref[II.6.6.5]{6.6.5}) remains valid when $K'$ is an
\emph{arbitrary} extension of $K$.

The rest of this subsection is devoted to the proof of the criterion
(\hyperref[II.6.6.11]{6.6.11}), which is a rather technical refinement of
(\hyperref[II.6.6.1]{6.6.1}); it may be omitted on a first reading.

\begin{lemma}[6.6.6]
\label{II.6.6.6}
Let $X$ be a reduced Noetherian prescheme and $\mathcal{E}$ a coherent
$\mathcal{O}_X$-module such that
$\mathcal{E}\otimes_{\mathcal{O}_X}\mathcal{R}(X)$ is a locally free
$\mathcal{R}(X)$-module of constant rank $n$. Then there exist
a reduced Noetherian prescheme $Z$ and a finite birational morphism
$h:Z\to X$ having the following property: the morphisms of sheaves
of sets
$\sigma_i:\mathcal{H}om_{\mathcal{O}_X}(\mathcal{E},\mathcal{E})
\to\mathcal{R}(X)$ ($1\leq i\leq n$)
(cf. \hyperref[II.6.4.10]{6.4.10}) map
$\mathcal{H}om_{\mathcal{O}_X}(\mathcal{E},\mathcal{E})$ into the
coherent $\mathcal{O}_X$-algebra $h_*(\mathcal{O}_Z)$.
\end{lemma}

Indeed, consider an affine open subset $U$ of $X$, with ring $A(U)=A$; let
$E=\Gamma(U,\mathcal{E})$, and let $C_U$ be the subalgebra of $R(U)$ generated
by the $\sigma_i(u)$ as $u$ runs through $\operatorname{Hom}_A(E,E)$; we
have seen (\hyperref[II.6.4.7.1]{6.4.7.1}) that this $A$-algebra is of
\emph{finite rank}. Moreover, it is clear that the formation of the algebras
$C_U$ commutes with the operations of restriction from an affine open subset
$U$ to an affine open subset $U'\subset U$. We have thus defined a
\emph{finite} sub-$\mathcal{O}_X$-algebra $\mathcal{C}$ of $\mathcal{R}(X)$
such that $\Gamma(U,\mathcal{C})=C_U$ for every affine open subset $U$ of
$X$. We

\oldpage[II]{133}

take $Z=\operatorname{Spec}(\mathcal{C})$, and for $h$ the structural
morphism, which is therefore finite (\hyperref[II.6.1.2]{6.1.2}); since
$\mathcal{C}$ is reduced, $Z$ is a reduced Noetherian prescheme
(\hyperref[II.1.3.8]{1.3.8}). Finally, the total ring of fractions of $C_U$
is $R(U)$ by definition, and since $C_U$ is contained in the integral
closure of $A(U)$ in $R(U)$, there is a bijective correspondence between the
minimal prime ideals of $A(U)$ and the minimal prime ideals of $C_U$
(\cite[t.~I, p.~259]{I-13}), which proves that $h$ is birational and
completes the proof.

\begin{corollary}[6.6.7]
\label{II.6.6.7}
Under the hypotheses of (\hyperref[II.6.6.6]{6.6.6}), let $W$ be an open
subset of $X$ such that, for every $x\in W$, either $X$ is normal at the
point $x$, or $\mathcal{E}_x$ is a free $\mathcal{O}_x$-module. Then we may
suppose $h$ defined in such a way that the restriction of $h$ to $h^{-1}(W)$
is an isomorphism of $h^{-1}(W)$ onto $W$.
\end{corollary}

Indeed, either hypothesis implies that, if $U\subset W$ is an affine open
subset, we have, with the notation of
(\hyperref[II.6.6.6]{6.6.6}), $(\sigma_i(u))_x\in A_x$ for every
$x\in U$ (\hyperref[II.6.4.3]{6.4.3}), hence $\sigma_i(u)\in A$, and the
conclusion follows from the definition of $h$ given in
(\hyperref[II.6.6.6]{6.6.6}).

\begin{env}[6.6.8]
\label{II.6.6.8}
Let $X$ be a reduced Noetherian prescheme and $g:X'\to X$ a finite
surjective morphism, so that $\mathcal{B}=g_*(\mathcal{O}_{X'})$ is a
\emph{coherent} $\mathcal{O}_X$-algebra; suppose moreover that
$\mathcal{B}\otimes_{\mathcal{O}_X}\mathcal{R}(X)$ is an
$\mathcal{R}(X)$-module that is \emph{locally free of constant rank $n$}. We
can then apply Lemma (\hyperref[II.6.6.6]{6.6.6}) taking
$\mathcal{E}=\mathcal{B}$, whence, with the notation of
(\hyperref[II.6.6.6]{6.6.6}), a homomorphism of sheaves of multiplicative
monoids
$\sigma_n:\mathcal{H}om_{\mathcal{O}_X}(\mathcal{B},\mathcal{B})
\to h_*(\mathcal{O}_Z)$, and, by composing this homomorphism with
the canonical homomorphism
$\mathcal{B}\to
\mathcal{H}om_{\mathcal{O}_X}(\mathcal{B},\mathcal{B})$
(\hyperref[II.6.5.1]{6.5.1}), we therefore obtain a homomorphism of
sheaves of multiplicative monoids:
\[
\label{II.6.6.8.1}
N':\mathcal{B}=g_*(\mathcal{O}_{X'})\longrightarrow
h_*(\mathcal{O}_Z)=\mathcal{C}.
\tag{6.6.8.1}
\]

This being so, for every invertible $\mathcal{O}_{X'}$-module $\mathcal{L}'$,
$g_*(\mathcal{L}')$ is an invertible $\mathcal{B}$-module
(\hyperref[II.6.1.12]{6.1.12}), and the method of
(\hyperref[II.6.5.2]{6.5.2}) makes it possible to associate functorially to
$\mathcal{L}'$ an invertible $\mathcal{C}$-module, which we will denote by
$N'(g_*(\mathcal{L}'))$.
\end{env}

\begin{lemma}[6.6.9]
\label{II.6.6.9}
Let $X$ be a reduced Noetherian prescheme and $g:X'\to X$ a finite
surjective morphism such that
$g_*(\mathcal{O}_{X'})\otimes_{\mathcal{O}_X}\mathcal{R}(X)$ is a
locally free $\mathcal{R}(X)$-module of constant rank. Then there exist
a reduced Noetherian prescheme $Z$ and a finite birational morphism
$h:Z\to X$ having the following property: for every ample
$\mathcal{O}_{X'}$-module $\mathcal{L}'$, the invertible
$\mathcal{O}_Z$-module $\mathcal{M}$ such that
$h_*(\mathcal{M})=N'(g_*(\mathcal{L}'))$ (notation of
\hyperref[II.6.6.8]{6.6.8}) is ample.
\end{lemma}

Suppose first that the homomorphism
$\mathcal{B}\to
\mathcal{B}\otimes_{\mathcal{O}_X}\mathcal{R}(X)$ is \emph{injective}.
Define $Z$ and $h$ as in (\hyperref[II.6.6.6]{6.6.6}) (with
$\mathcal{E}=g_*(\mathcal{O}_{X'})$). Let $z\in Z$; we must show that there
exist an integer $m>0$ and a section $t$ of $\mathcal{M}^{\otimes m}$
over $Z$ such that $Z_t$ is an affine open subset containing $z$
(\hyperref[II.4.5.2]{4.5.2}). Let $x=h(z)$, and let $U$ be an affine open
subset of $X$ containing $x$; $h^{-1}(U)$ is then an affine open
neighbourhood of $z$, and it will suffice to find $t$ such that
$z\in Z_t\subset h^{-1}(U)$, since $Z_t$ will then necessarily be affine
(\hyperref[II.5.5.8]{5.5.8}). By hypothesis there exist an integer $n>0$
and a section $s'$ of ${\mathcal{L}'}^{\otimes n}$ over $X'$ such that
we have
\[
\label{II.6.6.9.1}
g^{-1}(x)\subset X'_{s'}\subset g^{-1}(U)
\tag{6.6.9.1}
\]
by virtue of (\hyperref[II.4.5.4]{4.5.4}). By definition, $s'$ is also a
section of $g_*({\mathcal{L}'}^{\otimes n})$ over $X$, and there
corresponds to it, as in (\hyperref[II.6.5.2]{6.5.2}), a section
$s=N'(s')$ of $N'(g_*({\mathcal{L}'}^{\otimes n}))$ over

\oldpage[II]{134}

$X$. We will see that, if $t$ is the section $s$ considered as a section
of $\mathcal{M}^{\otimes n}$ over $Z$, then $t$ answers the question.
Indeed, set
\[
\label{II.6.6.9.2}
V=X-g(X'-X'_{s'})
\tag{6.6.9.2}
\]
which is an open subset of $X$ containing $x$ and contained in $U$, by
virtue of (\hyperref[II.6.6.9.1]{6.6.9.1}) and
(\hyperref[II.6.1.10]{6.1.10}). We will show that we have
\[
\label{II.6.6.9.3}
h^{-1}(V)\subset Z_t\subset h^{-1}(U),
\tag{6.6.9.3}
\]
which will complete the proof. It amounts to the same to say that the set
$T$ of $y\in X$ such that $s_y$ is invertible contains $V$ and is contained
in $U$. For this, consider first an affine open subset $W$ contained in
$V$; $g^{-1}(W)$ is an affine open subset of $X'$, and, by virtue of
(\hyperref[II.6.6.9.2]{6.6.9.2}), $s'_{y'}$ is invertible for every
$y'\in g^{-1}(W)$; by virtue of the hypotheses on $X$ and $\mathcal{B}$, we
can apply the results of (\hyperref[II.6.4.7]{6.4.7}), and we see that,
if $y=g(y')$, $s_y$ is invertible; in other words $V\subset T$. On the other
hand, it also follows from (\hyperref[II.6.4.7]{6.4.7}) that, conversely,
if $s_y$ is invertible, then so is $s'_{y'}$, which implies
$y'\in g^{-1}(U)$ by (\hyperref[II.6.6.9.1]{6.6.9.1}), and consequently
$y\in U$; whence $T\subset U$ in this case.

We pass from this to the general case by the same reasoning as in
(\hyperref[II.6.6.2]{6.6.2}), replacing $X'$ by an $X''$ such that
\[
g'_*(\mathcal{O}_{X''})\longrightarrow
g'_*(\mathcal{O}_{X''})\otimes_{\mathcal{O}_X}\mathcal{R}(X)
\]
is injective, and $\mathcal{L}'$ by an ample $\mathcal{O}_{X''}$-module
$\mathcal{L}''$ such that
$N'(g_*(\mathcal{L}'))=N'(g'_*(\mathcal{L}''))$. Lemma
(\hyperref[II.6.6.9]{6.6.9}) is therefore proved in all cases (with a
suitable choice of $h$).

\begin{corollary}[6.6.10]
\label{II.6.6.10}
Suppose that the hypotheses of
(\hyperref[II.6.6.9]{6.6.9}) are satisfied; for every invertible
$\mathcal{O}_X$-module $\mathcal{L}$ such that $g^*(\mathcal{L})$ is ample,
$h^*(\mathcal{L})$ is ample.
\end{corollary}

Indeed, if we set $\mathcal{L}'=g^*(\mathcal{L})$, we then have
$g_*(\mathcal{L}')=\mathcal{L}\otimes_{\mathcal{O}_X}\mathcal{B}$
(\hyperref[0.5.4.10]{0, 5.4.10}), hence
\[
N'(g_*(\mathcal{L}'))=
(\mathcal{L}\otimes_{\mathcal{O}_X}\mathcal{C})^{\otimes n}
\]
(by the same reasoning as for
\hyperref[II.6.5.2.4]{6.5.2.4}). We conclude that we have
$\mathcal{M}=(h^*(\mathcal{L}))^{\otimes n}$, and since $\mathcal{M}$ is
ample, so is $h^*(\mathcal{L})$
(\hyperref[II.4.5.6]{4.5.6}).

\begin{proposition}[6.6.11]
\label{II.6.6.11}
Let $Y$ be an affine scheme, $X$ a reduced Noetherian prescheme,
$f:X\to Y$ a quasi-compact morphism, and $g:X'\to X$ a finite surjective
morphism. Let $W$ be an open subset of $X$ such that, for every $x\in W$,
either $X$ is normal at the point $x$, or there exists an open neighbourhood
$T\subset W$ of $x$ such that $(g_*(\mathcal{O}_{X'}))|T$ is a locally free
$(\mathcal{O}_X|T)$-module. Then there exist a
reduced $Y$-prescheme $Z$ and a finite and birational $Y$-morphism
$h:Z\to X$ such that the restriction of $h$ to $h^{-1}(W)$ is an isomorphism
$h^{-1}(W)\xrightarrow{\sim}W$, and which has the following property: for
every invertible $\mathcal{O}_X$-module $\mathcal{L}$ such that
$g^*(\mathcal{L})$ is ample relative to $f\circ g$, $h^*(\mathcal{L})$
is ample relative to $f\circ h$.
\end{proposition}

Since $Y$ is affine, $g^*(\mathcal{L})$ is ample, and it is then a matter of
proving that, for a suitable choice of $h$, $h^*(\mathcal{L})$ is ample
(\hyperref[II.4.6.6]{4.6.6}). We will see that we can replace $g$
by a finite surjective morphism $g':X''\to X$ such that
${g'}^*(\mathcal{L})$ is ample and
$g'_*(\mathcal{O}_{X''})\otimes_{\mathcal{O}_X}\mathcal{R}(X)$ is an
$\mathcal{R}(X)$-module \emph{locally free of constant rank}; we will
then be reduced to the conditions of (\hyperref[II.6.6.10]{6.6.10}), and the
proposition will be proved.

For this, let $\mathcal{B}=g_*(\mathcal{O}_{X'})$; denote by $X_i$
($1\leq i\leq n$) the reduced closed sub-preschemes of $X$ whose
underlying spaces are the irreducible components of $X$
(\hyperref[I.5.2.1]{I, 5.2.1}); they are integral by hypothesis. Let
$X'_i$ be the closed sub-prescheme

\oldpage[II]{135}

$g^{-1}(X_i)$ of $X'$, and $g_i:X'_i\to X_i$ the morphism $g$ restricted to
$X'_i$, which is finite (\hyperref[II.6.1.5]{6.1.5, (iii)}) and surjective;
let $k_i$ be the \emph{rank} of the $\mathcal{O}_{X_i}$-algebra
$\mathcal{B}_i=(g_i)_*(\mathcal{O}_{X'_i})$. Since
$\mathcal{B}_i\otimes_{\mathcal{O}_{X_i}}\mathcal{R}(X_i)$ is a
constant presheaf (\hyperref[I.7.3.5]{I, 7.3.5}), the rank $k_i$ is
also the rank of the $(\mathcal{O}_X|U)$-algebra $\mathcal{B}|U$ for every
open subset $U$ of $X$ meeting \emph{only} the irreducible component
$X_i$. If $T$ is an open set in $X$ such that $\mathcal{B}|T$ is
isomorphic to $\mathcal{O}_X^m|T$, it follows from the preceding remark that
the numbers $k_i$ are equal to $m$ for all indices $i$ such that
$T\cap X_i\neq\varnothing$. Then let $U$ be the open subset of $X$ consisting of
the points in a neighbourhood of which $\mathcal{B}$ is a locally free
$\mathcal{O}_X$-module, and let $U_j$ ($1\leq j\leq s$) be its connected
components, which are open in $X$ and finite in number (since $U$ is
Noetherian); denote by $V_j$ the closed sub-prescheme of $X'$ that is the
\emph{closure} of the sub-prescheme induced on the open subset $g^{-1}(U_j)$
(\hyperref[I.9.5.11]{I, 9.5.11}). By the foregoing, for all the
indices $i$ such that $X_i\cap U_j\neq\varnothing$, the ranks $k_i$ are all
equal to one and the same integer $m_j$; note moreover that one and the same
$X_i$ cannot meet two $U_j$ with distinct indices. Let $i_\lambda$ be the
indices $i$ such that $X_i\cap U=\varnothing$. Consider the product $k$ of all the
$k_i$, set $n_i=k/k_i$, and let $X''$ be the prescheme defined as follows.
For each $j$ ($1\leq j\leq s$), we consider $k/m_j$ preschemes
isomorphic to $V_j$, and, for each $\lambda$, $k/k_{i_\lambda}$ preschemes
isomorphic to $X'_{i_\lambda}$; $X''$ is the \emph{sum} of all these
preschemes. We define a morphism $g'':X''\to X'$ that reduces to the
canonical injection on each of the summands of $X''$; it is clear that
$g''$ is a finite dominant morphism, hence surjective (since a finite
morphism is closed (\hyperref[II.6.1.10]{6.1.10})); set
$g'=g\circ g''$, which is a finite surjective morphism $X''\to X$; we have
${g'}^*(\mathcal{L})={g''}^*(g^*(\mathcal{L}))$, hence
${g'}^*(\mathcal{L})$ is an ample $\mathcal{O}_{X''}$-module
(\hyperref[II.5.1.12]{5.1.12}). It is then clear that, for this new
prescheme $X''$, the ranks defined like the $k_i$ for $X'$ are all
\emph{equal to $k$}; taking into account
(\hyperref[I.7.3.3]{I, 7.3.3}), we conclude immediately that, for every
affine open subset $T$ of $X$,
\[
\bigl(g'_*(\mathcal{O}_{X''})
\otimes_{\mathcal{O}_X}\mathcal{R}(X)\bigr)|T
\]
is an $(\mathcal{R}(X)|T)$-module isomorphic to
$(\mathcal{R}(X)|T)^k$.

\begin{corollary}[6.6.12]
\label{II.6.6.12}
If, in the statement of (\hyperref[II.6.6.11]{6.6.11}), we have $W=X$, then,
for an invertible $\mathcal{O}_X$-module $\mathcal{L}$ to be ample
relative to $f$, it is necessary and sufficient that $g^*(\mathcal{L})$ be
ample relative to $f\circ g$.
\end{corollary}

\begin{remark}[6.6.13]
\label{II.6.6.13}
In Chapter III, we will see that if $Y$ is Noetherian, $f$ of finite type,
and the restriction of $f$ to the reduced closed sub-prescheme of $X$ having
$X-W$ as underlying space is \emph{proper}, then the conclusion of
(\hyperref[II.6.6.12]{6.6.12}) is still valid. But we will give in
Chapter V examples of algebraic schemes $X$ over a field $K$ (the
structural morphism $X\to\operatorname{Spec}(K)$ not being proper) whose
normalization $X'$ is quasi-affine, but which are not quasi-affine (so
that $\mathcal{O}_X$ is not ample, although $\mathcal{O}_{X'}$ is
(\hyperref[II.5.1.12]{5.1.12}) and the morphism $X'\to X$ is finite
and surjective (\hyperref[II.6.3.10]{6.3.10})). We will see in the next
subsection that this circumstance cannot occur when
``quasi-affine'' is replaced by ``affine''.
\end{remark}

\subsection{Chevalley's theorem.}
\label{subsection:II.6.7}

We will establish (with the help of Serre's criterion
(\hyperref[II.5.2.1]{5.2.1})) the following theorem, which was proved by
C.~Chevalley by other methods in the case of \emph{algebraic}
schemes.

\oldpage[II]{136}

\begin{theorem}[6.7.1]
\label{II.6.7.1}
Let $X$ be an affine scheme, $Y$ a Noetherian prescheme, and $f:X\to Y$ a
finite surjective morphism. Then $Y$ is an affine scheme.
\end{theorem}

It is clear that
$f_{\mathrm{red}}:X_{\mathrm{red}}\to Y_{\mathrm{red}}$ is finite
(\hyperref[II.6.1.5]{6.1.5, (vi)}); since $X_{\mathrm{red}}$ is an
affine scheme, and since it amounts to the same to say that $Y$ is affine or
that $Y_{\mathrm{red}}$ is affine, because $Y$ is Noetherian
(\hyperref[I.6.1.7]{I, 6.1.7}), we see that we may suppose $Y$
\emph{reduced}. For every closed subset $Y'$ of $Y$, there is then only one
reduced sub-prescheme of $Y$ having $Y'$ as underlying space
(\hyperref[I.5.1.2]{I, 5.1.2}); its inverse image $f^{-1}(Y')$,
canonically isomorphic to $X\times_Y Y'$
(\hyperref[I.4.4.1]{I, 4.4.1}), is affine, being a closed sub-prescheme of
$X$, and the restriction of $f$ to $f^{-1}(Y')$, which is identified with
$f\times_Y 1_{Y'}$, is a finite surjective morphism
(\hyperref[II.6.1.5]{6.1.5, (iii)}). By virtue of the principle of
Noetherian induction (\hyperref[0.2.2.2]{0, 2.2.2}), we can therefore
(taking into account
\hyperref[I.6.1.7]{I, 6.1.7}) reduce to proving the theorem under
the hypothesis that, for \emph{every closed subset $Y'\neq Y$}, every
closed sub-prescheme of $Y$ having $Y'$ as underlying space is affine.
We conclude from this that, \emph{for every coherent $\mathcal{O}_Y$-module
$\mathcal{F}$ whose (closed) support $Z$ is distinct from $Y$, we have
$\mathrm{H}^1(Y,\mathcal{F})=0$}. Indeed, there exists a closed sub-prescheme
$Y'$ of $Y$ having $Z$ as underlying space and such that, if $j:Y'\to Y$ is
the canonical injection, we have
$\mathcal{F}=j_*(j^*(\mathcal{F}))$
(\hyperref[I.9.3.5]{I, 9.3.5}); consequently
(\hyperref[II.5.2.3]{5.2.3}),
$\mathrm{H}^1(Y,\mathcal{F})
=\mathrm{H}^1(Y',j^*(\mathcal{F}))=0$ by
(\hyperref[I.5.1.9.2]{I, 5.1.9.2}).

Suppose first that $Y$ is not irreducible, and let $Y'$ be an irreducible
component of $Y$ and $Y''=Y-Y'$; we will also denote by
$Y'$ the reduced closed sub-prescheme of $Y$ having $Y'$ as underlying
space, and by $j$ the canonical injection $Y'\to Y$. Let $\mathcal{F}$
be a coherent $\mathcal{O}_Y$-module, and consider the canonical homomorphism
\[
\rho:\mathcal{F}\longrightarrow
\mathcal{F}'=j_*(j^*(\mathcal{F}))
\]
(\hyperref[0.4.4.3]{0, 4.4.3}); $\mathcal{F}'$ is a coherent
$\mathcal{O}_Y$-module by virtue of
(\hyperref[0.5.3.10]{0, 5.3.10}) and
(\hyperref[0.5.3.12]{0, 5.3.12}), since we have
$j_*(\mathcal{O}_{Y'})=\mathcal{O}_Y/\mathcal{J}$, denoting by
$\mathcal{J}$ the sheaf of ideals of $\mathcal{O}_Y$ defining the
sub-prescheme $Y'$. Consequently $\mathcal{G}=\operatorname{Ker}\rho$ and
$\mathcal{K}=\operatorname{Im}\rho$ are also coherent $\mathcal{O}_Y$-modules
(\hyperref[0.5.3.4]{0, 5.3.4}); now, by definition, the stalk
$\mathcal{F}'_y$ of $\mathcal{F}'$ at the generic point $y$ of $Y'$ is equal
to $\mathcal{F}_y$, because $y$ is an interior point of $Y'$, and we therefore
have $\mathcal{J}_y=0$, since $Y$ is reduced. We conclude that $y$ is not
contained in the (closed) support of $\mathcal{G}$; on the other hand, the
support of $\mathcal{F}'$ (and \emph{a fortiori} that of $\mathcal{K}$) is
contained in $Y'$; in other words, the supports of $\mathcal{G}$ and of
$\mathcal{K}$ are \emph{distinct from $Y$}. We deduce that
$\mathrm{H}^1(Y,\mathcal{G})=\mathrm{H}^1(Y,\mathcal{K})=0$, and the exact
cohomology sequence applied to the exact sequence
$0\to\mathcal{G}\to\mathcal{F}\to\mathcal{K}\to0$ shows that
$\mathrm{H}^1(Y,\mathcal{F})=0$. We then conclude by Serre's criterion
(\hyperref[II.5.2.1]{5.2.1}).

Suppose therefore that $Y$ is irreducible, and consequently \emph{integral}.
We may also suppose that $X$ is \emph{integral}: indeed, if we denote by $X_i$
the reduced closed sub-preschemes of $X$ having as underlying spaces the
irreducible components of $X$ (\hyperref[I.5.2.1]{I, 5.2.1}) and by
$f_i$ the restriction of $f$ to $X_i$, at least one of the $f_i$ is dominant,
and, since it is a finite morphism (\hyperref[II.6.1.5]{6.1.5}), it is
surjective (\hyperref[II.6.1.10]{6.1.10}); since, on the other hand, $X_i$ is
an affine scheme, we see that we can replace $X$ by $X_i$ in the statement.

\begin{lemma}[6.7.1.1]
\label{II.6.7.1.1}
Let $X$, $Y$ be two integral Noetherian preschemes, $x$ (resp. $y$) the
generic point of $X$ (resp. $Y$), and $f:X\to Y$ a finite surjective morphism.
Let $\mathcal{L}$ be an invertible $\mathcal{O}_X$-module such that there
exist an affine open neighbourhood $U$ of $y$ and a section
$g\in\Gamma(X,\mathcal{L})$ for which

\oldpage[II]{137}

$x\in X_g\subset f^{-1}(U)$. Then there exist two integers $m>0$, $n>0$, a
homomorphism
$u:\mathcal{O}_Y^m\to f_*(\mathcal{L}^{\otimes n})$ and an open neighbourhood
$V$ of $y$ such that the restriction $u|V$ is an isomorphism of
$\mathcal{O}_Y^m|V$ onto $f_*(\mathcal{L}^{\otimes n})|V$.
\end{lemma}

Let $C$ be the (integral) ring of $U$ and $k=\mathcal{O}_y$ its field of
fractions: since $f$ is finite, $U'=f^{-1}(U)$ is affine
(\hyperref[II.1.3.2]{1.3.2}); let $D$ be its (integral) ring, with field of
fractions $K=\mathcal{O}_x$; by hypothesis $D$ is a $C$-module of finite
type (\hyperref[II.6.1.4]{6.1.4}), hence $K$ is an extension of finite rank
of $k$. The fibre
\[
f^{-1}(y)=X\times_Y\operatorname{Spec}(k(y))
=X\times_Y\operatorname{Spec}(\mathcal{O}_y)
\]
is identified with $\operatorname{Spec}(K)$
(\hyperref[I.3.6.5]{I, 3.6.6}); let $s_i$ ($1\leq i\leq m$) be
elements of $D$ forming a basis of $K$ over $k$. There exists $n>0$ such that
the sections $(s_i|X_g)g^{\otimes n}$ of $\mathcal{L}^{\otimes n}$ over
$X_g$ extend to sections $b_i$ ($1\leq i\leq m$) of
$\mathcal{L}^{\otimes n}$ over $X$
(\hyperref[I.9.3.1]{I, 9.3.1}). By definition, the $b_i$ are also
sections of $f_*(\mathcal{L}^{\otimes n})$ over $Y$, and they therefore
define a homomorphism
$u:\mathcal{O}_Y^m\to f_*(\mathcal{L}^{\otimes n})$
(\hyperref[0.5.1.1]{0, 5.1.1}); we will see that $u$ answers the
question. We have $\mathcal{L}^{\otimes n}|U'=\widetilde{M}$, where $M$ is a
$D$-module of finite type; hence, if $\varphi$ is the injection $C\to D$
corresponding to the morphism $f^{-1}(U)\to U$ that is the restriction of $f$,
$M_{[\varphi]}$ is a $C$-module of finite type; consequently
\[
f_*(\mathcal{L}^{\otimes n})|U=(M_{[\varphi]})^\sim
\]
(\hyperref[I.1.6.3]{I, 1.6.3}) is coherent, and, since $U$ is an arbitrary
affine open subset of $Y$, $f_*(\mathcal{L}^{\otimes n})$ is coherent;
moreover $u|U=\widetilde{\theta}$, where $\theta$ is a $C$-homomorphism
$C^m\to M_{[\varphi]}$, and $u_y=\theta_y$ is the homomorphism
\[
\theta\otimes1:K^m=C^m\otimes K\longrightarrow
M_{[\varphi]}\otimes K,
\]
and the latter is by definition an \emph{isomorphism}, because the $(b_i)_x$
form a basis of $M_{[\varphi]}\otimes K$ over $k$, $g_x$ being by
hypothesis a generator of $(\mathcal{L}^{\otimes n})_x$. We conclude that
the supports of the $\mathcal{O}_Y$-modules $\operatorname{Ker}u$ and
$\operatorname{Coker}u$ do not contain $y$; since these
$\mathcal{O}_Y$-modules are coherent
(\hyperref[0.5.3.3]{0, 5.3.3}), their supports are closed
(\hyperref[0.5.2.2]{0, 5.2.2}), whence the lemma.

This being so, the hypotheses of Lemma
(\hyperref[II.6.7.1.1]{6.7.1.1}) are satisfied in the case we are
considering by taking $\mathcal{L}=\mathcal{O}_X$, since $X$ is affine
(\hyperref[I.1.1.10]{I, 1.1.10}); we set
$\mathcal{A}=\mathcal{O}_Y$, $\mathcal{B}=f_*(\mathcal{O}_X)$. By virtue of
Serre's criterion (\hyperref[II.5.2.1.1]{5.2.1.1}), it suffices to prove
that, for every coherent $\mathcal{O}_Y$-module $\mathcal{F}$, we have
$\mathrm{H}^1(Y,\mathcal{F})=0$; it even suffices to prove this when
$\mathcal{F}\subset\mathcal{O}_Y$, which implies that $\mathcal{F}$ is
torsion-free, since $Y$ is integral; in fact we will show that
$\mathrm{H}^1(Y,\mathcal{F})=0$ for \emph{every} coherent
$\mathcal{O}_Y$-module $\mathcal{F}$ that is \emph{torsion-free}. Now the
homomorphism $u:\mathcal{A}^m\to\mathcal{B}$ defines a homomorphism
\[
v:\mathcal{G}
=\mathcal{H}om_{\mathcal{A}}(\mathcal{B},\mathcal{F})
\longrightarrow
\mathcal{H}om_{\mathcal{A}}(\mathcal{A}^m,\mathcal{F})
=\mathcal{F}^m.
\]
Let us first show that $v$ is \emph{injective}: by hypothesis,
$\mathcal{T}=\operatorname{Coker}u$ has a support not meeting $V$, hence
is a torsion $\mathcal{O}_Y$-module
(\hyperref[I.7.4.6]{I, 7.4.6}); the exact sequence
\[
\mathcal{A}^m\longrightarrow\mathcal{B}\longrightarrow\mathcal{T}
\longrightarrow0
\]
gives, by left exactness of the functor
$\mathcal{H}om_{\mathcal{A}}$, the exact sequence
\[
0\longrightarrow
\mathcal{H}om_{\mathcal{A}}(\mathcal{T},\mathcal{F})
\longrightarrow\mathcal{G}\xrightarrow{v}\mathcal{F}^m.
\]
But since $\mathcal{F}$ is torsion-free, we have
$\mathcal{H}om_{\mathcal{A}}(\mathcal{T},\mathcal{F})=0$
(\hyperref[0.5.2.6]{0, 5.2.6}), whence our assertion. We therefore have
the exact sequence
\[
0\longrightarrow\mathcal{G}\longrightarrow\mathcal{F}^m
\longrightarrow\operatorname{Coker}v\longrightarrow0
\]

\oldpage[II]{138}

where $\mathcal{G}$ and $\operatorname{Coker}v$ are coherent
$\mathcal{O}_Y$-modules
(\hyperref[0.5.3.4]{0, 5.3.4} and
\hyperref[0.5.3.5]{5.3.5}); by virtue of the exact cohomology sequence,
it will suffice to show that
$\mathrm{H}^1(Y,\mathcal{G})
=\mathrm{H}^1(Y,\operatorname{Coker}v)=0$ in order to deduce that
$\mathrm{H}^1(Y,\mathcal{F}^m)
=(\mathrm{H}^1(Y,\mathcal{F}))^m=0$, and consequently that
$\mathrm{H}^1(Y,\mathcal{F})=0$. Now the restriction $v|V$ is an isomorphism,
hence the support of $\operatorname{Coker}v$ is distinct from $Y$, whence
$\mathrm{H}^1(Y,\operatorname{Coker}v)=0$ by the hypothesis made at the
beginning. On the other hand, $\mathcal{G}$ is a coherent $\mathcal{B}$-module
(\hyperref[I.9.6.4]{I, 9.6.4}); since $X$ is affine over $Y$, there
exists a quasi-coherent $\mathcal{O}_X$-module $\mathcal{K}$ such that
$\mathcal{G}$ is isomorphic to $f_*(\mathcal{K})$
(\hyperref[II.1.4.3]{1.4.3}); since $X$ is affine, we have
$\mathrm{H}^1(X,\mathcal{K})=0$
(\hyperref[I.5.1.9.2]{I, 5.1.9.2}), hence also
$\mathrm{H}^1(Y,\mathcal{G})=0$ by
(\hyperref[II.5.2.3]{5.2.3}), which completes the proof of Theorem
(\hyperref[II.6.7.1]{6.7.1}).

\begin{corollary}[6.7.2]
\label{II.6.7.2}
Let $X$ be a Noetherian prescheme and $(X_i)_{1\leq i\leq n}$ a finite
covering of the space $X$ consisting of closed sets. For $X$ to be affine, it
is necessary and sufficient that, for each $i$, there exist a closed
sub-prescheme of $X$ having $X_i$ as underlying space and which is affine.
\end{corollary}

Indeed, if this is so, let $X'$ be the scheme that is the \emph{sum} of the
$X_i$; it is clear that $X'$ is affine, and we define a surjective morphism
$f:X'\to X$ by taking as the restriction of $f$ to $X_i$ the canonical
injection. Everything comes down to verifying that $f$ is \emph{finite}, by
virtue of (\hyperref[II.6.7.1]{6.7.1}), and this was seen in
(\hyperref[II.6.1.5]{6.1.5}).

\begin{corollary}[6.7.3]
\label{II.6.7.3}
For a Noetherian prescheme $X$ to be affine, it is necessary and sufficient
that the reduced closed sub-preschemes having as underlying spaces the
irreducible components of $X$ be affine.
\end{corollary}

